%% QDCVR platform — Knowledge-Based Systems submission (elsarticle)
%% v2 restructure 2026-10-12: content-retrieval-based KB *management system* framing.
%% Compile: pdflatex main2 && bibtex main2 && pdflatex main2 && pdflatex main2
%% AUTHOR ACTION: replace author placeholders, ORCID, and funding fields.
\documentclass[preprint,12pt]{elsarticle}

\usepackage[T1]{fontenc}
\usepackage[utf8]{inputenc}
\usepackage{booktabs}
\usepackage{array}
\usepackage{graphicx}
\usepackage{amsmath}
\usepackage{amssymb}
\usepackage{xcolor}
\usepackage{url}
\usepackage[colorlinks=true,linkcolor=blue,citecolor=blue,urlcolor=blue]{hyperref}
\usepackage{caption}
\captionsetup{font=small,labelfont=bf}
\newcommand{\yes}{\checkmark}
\newcommand{\no}{--}

\graphicspath{{figs/}}
\setlength{\emergencystretch}{1.5em}

\journal{Knowledge-Based Systems}

\begin{document}

\begin{frontmatter}

\title{Design of a content-retrieval-based knowledge-base management system
for trustworthy enterprise document intelligence}

\author[inst1]{Author Name\corref{cor1}}
\ead{author@example.com}
\author[inst1]{Co-author Name}
\cortext[cor1]{Corresponding author.}
\affiliation[inst1]{organization={Affiliation Department, University,
City, Country}, city={City}, country={Country}}

\begin{abstract}
Enterprises accumulate process regulations, failure analyses, and operational
lessons as unstructured documents spread across shared drives, and answering
one engineering question requires locating, reading and citing the right
document. Retrieval-augmented generation addresses the locating step, but its
dominant ranking signal measures topic relatedness rather than evidence
presence. We present QDCVR, a self-hosted knowledge-base management system that
governs content across all three actions. Six functional subsystems---a
quality-gated parsing and ingestion pipeline, content-routed organization,
life-cycle management with three-source consistency audits, a document graph,
an operational-knowledge store with credibility tiers, and multi-lane
content-verified retrieval---expose 95 Model Context Protocol tools over one
storage layer, and every subsystem terminates at an explicit contract rather
than an implicit prompt. We evaluate the system on a deployed industrial
corpus, a 322-document organization campaign, a 100-paper acquisition
campaign, a cross-lingual corpus, and functional and non-functional campaigns.
Twenty-one of 100 submitted papers were never retrievable from the arXiv
corpus, and the accounting gap between submission, storage and index coverage
is reported rather than hidden. The verification gate exposes a same-domain
score-inflation failure consistent with known LLM-judge miscalibration:
non-gold scores are indistinguishable from gold (Mann--Whitney $p = 0.26$),
judged re-ranking lowers Hit@5 from 0.94 to 0.44, while reusing the same
verdicts as a filter restores 0.94 (sign test $p = 0.008$). Functional
campaigns confirm ingestion, deduplication, organization, graph construction,
distillation and abstention end to end; engine choice moves monetary cost by
one to two orders of magnitude at the same quality tier. Verifier calibration
is the open problem, and honesty needs an enforcement layer.
\end{abstract}

\begin{keyword}
Knowledge-base management system \sep Document ingestion and parsing \sep
Knowledge organization and lifecycle \sep Content-verified retrieval \sep
Trustworthy question answering \sep Model Context Protocol
\end{keyword}

\end{frontmatter}

%% ============================================================ body (v2 order)
\section{Introduction}
\label{sec:intro}

Enterprises accumulate process regulations, failure analyses, and operational
lessons as unstructured documents spread across shared drives. Answering a
concrete engineering question --- \emph{why did edge delamination appear in an
electroluminescence inspection of a laminated photovoltaic module?} ---
requires locating the one report among thousands that records the decision,
reading it, and citing it. Retrieval-augmented generation (RAG) addresses the
location step by embedding documents and the query into a shared vector space
\cite{rag-lewis,dpr}, and agentic stacks let a large language model orchestrate
retrieval tools iteratively \cite{mcp,singh-agentic-rag}. Yet the dominant
ranking signal, dense cosine similarity, measures \emph{topic relatedness}
rather than \emph{evidence presence}. A document that merely discusses similar
materials can outrank the one document that contains the answer, and a fluent
generator can synthesize an answer from the wrong source
\cite{siren,gao-ragsurvey}.

Three research communities attack the three actions this task requires, and
each leaves the others' problems intact. Retrieval research optimizes
\emph{finding}: sparse, dense and hybrid rankers, and more recently graph-based
and hierarchical indexes, are evaluated by whether a relevant passage appears
in a top-$k$ list. Knowledge management addresses \emph{organizing}:
classification schemes, controlled vocabularies and lessons-learned systems
govern how an organization's documents are arranged and reused, but they do not
close the loop to an answer that a machine must justify. Document AI addresses
\emph{reading}: layout-aware parsing recovers text, tables and formulas from
scanned pages \cite{mineru}, but a parsed document still has to be placed,
indexed, governed and trusted. In a deployment these three concerns are not
independent: a parse failure that is not treated as a failure becomes a
silently missing document; a catalog that is not governed makes candidate
generation arbitrary; and a ranking score that is treated as evidence produces
confident citations to documents that were never read.

This paper takes the position that the object all three actions should govern
is \emph{content}, and that a knowledge-base management system should express
that position as enforceable contracts rather than as good practice. Three
commitments follow, and each one is implemented and measured here. First,
\emph{content decides placement}: library membership, tags, descriptions and
graph edges are derived from what a document says, not from the folder it
arrived in, and the resulting catalog is the system's first retrieval
instrument. Second, \emph{content decides admissibility}: no similarity score,
however high, promotes a candidate into an answer, because the candidate's
fetched content must be judged to contain direct evidence first, and the
platform fetches that content itself. Third, \emph{absence of evidence is an
answer}: when verification eliminates every candidate, the system returns an
explicit not-found report rather than the closest text. The consequence of the
three commitments is that the system's claims are auditable: every placement,
verdict, cut and lifecycle operation leaves a record, and the collection can
be re-verified afterwards.

We built QDCVR (v2.3), a self-hosted knowledge-base management system that
implements this position end to end, and we evaluate it on a deployed
industrial corpus and on acquisition and functional campaigns. The system
spans six functional subsystems --- a quality-gated parsing and ingestion
pipeline, knowledge organization with content routing, tags and a document
graph, life-cycle management with three-source consistency audits, an
operational-knowledge store with credibility tiers, a multi-lane
content-verified retrieval layer, and a tool interface that exposes the whole
platform as 95 Model Context Protocol tools \cite{mcp} --- and it runs on a
single workstation. Two results frame the paper. On the acquisition side, the
pipeline is where trust is won or lost: of 100 papers submitted to a
100-document campaign, 21 failed to parse, and the design decision that
matters is that a failed parse is reported rather than stored as an empty
document (Section~\ref{sec:e1}). On the retrieval side, the
content-verification gate that the system is built around turns out to be a
better filter than a ranker: reusing the \emph{same} verdicts as a re-ranking
signal lowers Hit@5 from 0.94 to 0.44, while using them as a binary filter
restores 0.94 (sign test $p = 0.008$), and the underlying cause is a
same-domain score inflation in which non-gold candidates are statistically
indistinguishable from gold ones (Mann--Whitney $p = 0.26$)
(Section~\ref{sec:eval-verify}). We report both outcomes, including the one
that is negative for our own design.

\noindent The contributions are:
\begin{enumerate}
\item \textbf{A content-governed knowledge-base management architecture}
(Sections~\ref{sec:system}--\ref{sec:implementation}): requirements derived
from deployment, six design principles enforced in code, a four-layer
architecture, and an explicit data model whose contracts --- parse gates,
placement rules, three-source consistency, promotion conditions, task
semantics --- are stated as obligations rather than conventions.
\item \textbf{A quality-gated acquisition pipeline with a measured failure
budget} (Section~\ref{sec:acquisition}, E1): multi-format parsing with OCR,
exact and near-duplicate detection, structure-aware splitting of oversized
documents with part addressing, content-derived descriptions and tags, and an
index post-condition; evaluated on a 100-document campaign (21\% parse
failure) and on a cross-lingual corpus.
\item \textbf{Content-routed organization and life-cycle management}
(Section~\ref{sec:organization}, E2): placement by content, catalog and
metadata registries with progressive disclosure, tag governance, duplicate
hygiene, a document graph with an explicitly restricted role, an experience
store with credibility tiers, and lifecycle operations certified by a
read-only audit (322 documents in five libraries).
\item \textbf{Multi-channel content-verified retrieval} (Section~\ref{sec:lanes},
E4--E5): three lanes over one tool surface under enforced budgets, a
verification gate served by a resident judging worker, and an honest
abstention path --- compared against seven baseline pipelines and against a
bare agent and a pure-dense controller.
\item \textbf{A verifier-failure analysis and its role decomposition}
(Section~\ref{sec:eval-verify}, E4, E6): same-domain inflation quantified
inside a deployed pipeline, the filter/re-ranker asymmetry, a threshold that
never abstains out of corpus, and robustness of the failure to paraphrase, to
a second verification signal and to a second corpus.
\item \textbf{An operational account of the system} (Sections~\ref{sec:impl}
and~\ref{sec:discussion}, E7): resident-worker and subprocess engineering,
engine-agnostic harness integration with cost measured across three engines,
a catalogue of failure modes with their detection symptoms, and the
non-functional behaviour observed under concurrent writes.
\end{enumerate}

A preliminary demo describing the platform's organization and retrieval layers
was published as a CIKM~2026 demo \cite{qdcvr-demo}. This paper extends it with
four components that the demo does not contain: the acquisition pipeline with
its measured parse-failure budget and cross-lingual evaluation; life-cycle
management and integrity auditing as a designed subsystem rather than a
maintenance procedure; the operational-knowledge store with its credibility
tiers and distillation loop; and the verifier-failure analysis with the
filter/re-ranker decomposition. The demo also reports no negative result, which
is the centrepiece here.

\subsection{Research questions and reading guide}
\label{sec:rq}

The paper is organized as a design-to-evidence arc around five research
questions, each answered by one evidence group and reported in one place
(Table~\ref{tab:ataglance}).

\begin{description}
\item[RQ1 (Enforceable contracts).] How can the three content commitments be
turned into system contracts that hold across retrieval strategies and
independent agent harnesses, rather than into prompt conventions?
\item[RQ2 (Acquisition and organization at scale).] Do quality-gated
acquisition and content-routed organization produce collections that are
usable and internally consistent when they grow, and what does the gate
reject?
\item[RQ3 (Content-verified retrieval).] Does staged, multi-channel retrieval
with content verification outperform one-shot vector top-$k$ and unverified
fusion --- and by which mechanism?
\item[RQ4 (Cost and benefit of verification).] What does the gate cost and
buy, and is its failure mode structural rather than an artefact of one
checkpoint, domain or language?
\item[RQ5 (Operational contract).] Where do honest abstention, latency, cost
and stability actually live in a running system?
\end{description}

\begin{table}[!htbp]
\centering
\footnotesize
\setlength{\tabcolsep}{3pt}
\caption{This paper at a glance: each research question, the evidence group
that answers it, and the headline finding.}
\label{tab:ataglance}
\begin{tabular}{p{0.8cm}p{3.7cm}p{2.9cm}p{4.9cm}}
\toprule
RQ & Question & Evidence & Headline finding \\
\midrule
RQ1 & Are the content commitments enforceable contracts? &
Architecture, interface, implementation
(\S\ref{sec:system}--\S\ref{sec:implementation}) & Yes: gates, placement
rules, lane whitelists and fail-closed verdicts make them properties of the
enforcement layer \\
RQ2 & Do gated acquisition and content routing scale? &
E1, E2 (\S\ref{sec:acquisition}, \S\ref{sec:organization}) & Partly: routing
and audits hold at 322 documents, while 21\% of submitted papers fail to
parse and must be rejected visibly \\
RQ3 & Does staged verified retrieval beat one-shot top-$k$? &
E4 (\S\ref{sec:eval-retrieval}) & Yes on reachability, through staged
recall plus content adjudication; two-stage retrieval is an accelerator, not
the contribution \\
RQ4 & What does the gate cost and buy, and is the failure structural? &
E4, E6 (\S\ref{sec:eval-verify}) & Costs $<$1\,s; as a re-ranker it loses
Hit@5 ($0.94\!\to\!0.44$), as a filter it is neutral ($0.94$,
$p=0.008$); the loss tracks confuser density, not checkpoint or domain \\
RQ5 & Where do honesty, latency and cost live? &
E5, E7 (\S\ref{sec:eval-agent}, \S\ref{sec:implementation}) & Honesty lives
in the reading and answering layer, never in the threshold; latency is
dominated by LLM turns, and engine choice moves cost by two orders of
magnitude \\
\bottomrule
\end{tabular}
\end{table}

The remainder of the paper follows the arc.
Section~\ref{sec:related} positions the work against knowledge management,
document parsing, retrieval, verification and agent protocols.
Sections~\ref{sec:system}--\ref{sec:organization} present the architecture,
the acquisition pipeline and the organization subsystem, and answer RQ1 and
RQ2. Section~\ref{sec:lanes} presents the retrieval lanes and the verification
gate, walking one real question end to end before the measurements begin, and
states what the retrieval layer claims (RQ3). Section~\ref{sec:implementation}
reports the engineering that makes the contracts enforceable.
Section~\ref{sec:eval} carries the seven evidence groups, including the
verifier's own failure analysis (RQ4) and the operational measurements (RQ5).
Sections~\ref{sec:discussion}--\ref{sec:conclusion} distil the lessons,
state the threats to validity, and conclude.
Table~\ref{tab:modules} maps every system function to its components and its
evidence; Table~\ref{tab:ataglance} maps every research question to its
evidence.

\section{Related work}
\label{sec:related}

\subsection{Knowledge bases, knowledge organization, and knowledge management}
Large public knowledge bases such as DBpedia and YAGO are built by extracting
structured facts from semi-structured sources and maintaining them under an
ontology \citep{auer-dbpedia-kb-2007,hoffart-yago2-kb-2013}. The enterprise
content management literature treats organizational content as a management
problem rather than an extraction problem: ECM integrates content, process, and
records perspectives \citep{paivarinta-ecm-2005}, and later work characterizes
the field's research agenda \citep{tyrvaainen-ecm-2006}. Organizational
knowledge creation describes how tacit experience becomes explicit, shareable
knowledge through socialization, externalization, combination, and
internalization, sustained by a shared context \citep{nonaka-seci-1994,
nonaka-seci-ba-2000}. Complementary strands study metadata quality in digital
repositories \citep{park-metadata-quality-2009}, the emergent vocabulary of
collaborative tagging \citep{golder-tagging-2006}, human-in-the-loop review of
machine-produced labels \citep{mosqueira-rey-hitl-2022}, and the inspectability
of automated decisions \citep{barredo-arrieta-xai-2020}.

This literature names the problems our design must answer --- inconsistent
metadata, uncontrolled vocabulary, unverified labels, experience that never
becomes reusable --- but it treats them as management and social problems and
rarely closes the loop to machine answering. That loop is the contribution
here: the content-derived metadata that organizes a base also constrains
retrieval, verification, and abstention
(Sections~\ref{sec:organization}--\ref{sec:gate}).

\subsection{Document parsing and acquisition pipelines}
Document AI supplies the components of an acquisition pipeline. PubLayNet
established large-scale layout analysis \citep{zhong-publaynet-layout-2019},
and unified text--image pre-training improved layout-aware understanding of
tables and figures \citep{huang-layoutlmv3-docai-2022}. MinerU provides an open
extraction pipeline for scientific and scanned documents and is the parser this
platform drives \citep{mineru}. Near-duplicate detection by locality-sensitive
hashing over shingles is the standard mechanism for redundancy control at crawl
scale \citep{manku-near-dup-2007}, and provenance has mature foundations in the
W3C PROV family \citep{missier-prov-2013} and in database provenance
\citep{cheney-provenance-2009}.

Parsing is usually presented as preprocessing. We treat it as a precondition:
an extraction error propagates into chunking, description, indexing, and every
later citation, and it is invisible in retrieval metrics. In our own campaign
21 of 100 submitted arXiv papers never reached a retrievable state
(Section~\ref{sec:acquisition}). Acquisition therefore carries its own quality
gate, duplicate check, index-verification step, and per-document provenance
record.

\subsection{Retrieval: sparse, dense, hybrid, and RAG}
Retrieval-augmented generation grounds generation in retrieved passages
\citep{rag-lewis}; dense retrievers learn a shared query--passage space
\citep{dpr}; multi-lingual, multi-granularity embedders such as BGE-M3 combine
dense and sparse signals \citep{bge-m3}. Lexical probabilistic ranking remains
a strong baseline \citep{bm25}, reciprocal-rank fusion merges ranked lists
without training \citep{rrf-cormack}, and ranking in the large-data regime is
benchmarked on MS~MARCO \citep{craswell-msmarco-2021}. Graded-relevance metrics
are the standard evaluation instrument \citep{jarvelin-cumulated-gain-2002}, and
surveys systematize the design space as retrieve-then-read pipelines
\citep{gao-ragsurvey}. The platform stores vectors in a persistent store with
graph-based approximate nearest-neighbour search \citep{chroma,hnsw}.

On our gold-annotated corpus the lexical baseline is not the weak part: BM25
reaches Hit@5 of 1.00. This paper does not compete on ranking quality. It treats
all of these signals as candidate generators whose output must still be
confirmed against content (Section~\ref{sec:lanes}).

\subsection{Structured and graph-based organization}
GraphRAG builds entity--relation graphs and community summaries for
query-focused summarization \citep{graphrag}; LightRAG reduces the idea to
dual-level retrieval \citep{lightrag}; RAPTOR recursively abstracts document
trees \citep{raptor}. Knowledge-graph work supplies the representational
vocabulary --- embedding-based representation and link prediction
\citep{wang-kg-embedding-2017}, and the acquisition, representation, and
application taxonomy \citep{ji-kg-survey-2022} --- with property-graph stores as
the common substrate \citep{neo4j}. Structured organization is also motivated as
a mitigation for ungrounded generation \citep{siren}.

Our graph layer is deliberately thinner than these systems: Document,
KnowledgeBase, and Tag nodes with membership, tag, and weighted relatedness
edges, built from metadata that acquisition already verified. It navigates and
expands candidates and never answers alone. The measured basis is that
keyword-limited graph search is the weakest standalone signal on our corpus
(Hit@5 0.56), so the platform exposes it as one channel among several
(Section~\ref{sec:lanes}).

\subsection{Content verification, correction, and abstention}
Self-RAG trains a generator to critique retrieved evidence and its own output
\citep{selfrag}; CRAG adds a corrective retriever that re-evaluates and
re-queries when retrieval quality is low \citep{crag}; FLARE interleaves
generation with forward-looking retrieval triggers \citep{flare};
\citet{kamath} formalize selective prediction under domain shift; and CyberBOT
applies ontology-grounded retrieval in a reliability-critical domain
\citep{cyberbot}. Verifiability measurement in generative search
\citep{liu-verifiability} and citation-capable generation \citep{gao-citations}
frame the problem this paper attacks: citations that are present but not
supportive.

\emph{LLM-as-a-judge.} The verification gate is a locally hosted judgement
model \citep{laya}, and the judging literature documents exactly the failure
modes our measurements expose. MT-Bench and Chatbot Arena established the
paradigm and its bias taxonomy \citep{zheng-llmjudge}; LLM relevance judges
agree imperfectly with human relevance \citep{sun-urara}; listwise LLM
re-rankers trade quality for instability at scale \citep{rankgpt}; evaluators
favour generations that resemble their own distribution
\citep{panickssery-selfpref}; and verifiers can be gamed by optimization
pressure \citep{lightman-verify}. Our contribution is not the discovery that
judges miscalibrate. It is the operational quantification of same-domain
inflation inside a deployed pipeline and the decomposition of which judge
\emph{role} --- filter, ordering, or threshold --- is safe under which pool
conditions (Sections~\ref{sec:inflation} and \ref{sec:gate}).

These approaches verify \emph{during generation} with the generator itself.
This platform externalizes verification into a dedicated model that sees only
(query, document text) pairs and is consumed as a tool: any lane, agent, or
external system can call it, and its verdicts fail closed. Verification thus
becomes an inspectable, reusable service rather than a property of one
generator's prompting.

\emph{Abstention.} Beyond selective prediction under domain shift
\citep{kamath}, recent work probes whether LLMs know what they do not know
\citep{yin-knowwhat} and orchestrates models to abstain on knowledge gaps
\citep{feng-abstain}. Section~\ref{sec:abstention} locates the abstention of
this platform not in calibrated uncertainty but in the reading-and-answer
structure --- a complementary, architectural answer to the same problem.

\subsection{Agentic retrieval and tool protocols}
ReAct established interleaved reasoning and acting over tools \citep{react-yao};
tool-learning extended it to self-supervised tool use \citep{toolformer}; the
Model Context Protocol standardizes tool servers as a machine interface
\citep{mcp}; agentic-RAG surveys systematize the design space and list
efficiency among its open challenges \citep{singh-agentic-rag}; and Modular RAG
frames retrieval systems as reconfigurable modules \citep{modular-rag}. DeepRead
shows that structure-aware navigation over long documents improves agentic
search \citep{deepread}, which matches the document-coordinate reading of the
librarian lane. The contribution here is a policy layer rather than a new
reasoning scheme: per-lane tool whitelists enforced by the server, hard call
budgets, and single-call server-side orchestration for the hybrid lane, aimed
at the turn count that dominates agentic latency
(Section~\ref{sec:composition}).

\subsection{Positioning}
Table~\ref{tab:positioning} compares representative systems along the
dimensions a knowledge-base management system has to satisfy.

\begin{table}[!htbp]
\centering
\small
\setlength{\tabcolsep}{4pt}
\caption{Positioning against representative systems. \emph{Parse}: a quality
gate that rejects unusable documents before they enter the base. \emph{Org.\ \&
lifecycle}: organization and governance of the base beyond a flat index
(catalog, metadata and tag control, duplicate control, move/merge/delete).
\emph{Experience}: operational lessons kept as first-class retrievable
knowledge. \emph{Verify}: every cited item is judged against the query before it
may support an answer. \emph{Abstain}: explicit not-found is terminal, not a
fallback to parametric or web knowledge. \emph{Lanes}: several retrieval
strategies behind one interface under enforced budgets. \emph{Protocol}: a
standardized agent tool protocol as the machine interface. \yes{} present;
\no{} not addressed by the cited work; \textit{n/r} not reported in the
accessible source; italic entries name the mechanism where a bare mark would
mislead. Marks denote capability, not measured quality.}
\label{tab:positioning}
\begin{tabular}{@{}p{5.6cm}ccccccc@{}}
\toprule
 & \rotatebox{90}{Parse} & \rotatebox{90}{Org.\ \& lifecycle} &
\rotatebox{90}{Experience} & \rotatebox{90}{Verify} &
\rotatebox{90}{Abstain} & \rotatebox{90}{Lanes} &
\rotatebox{90}{Protocol} \\
\midrule
Vanilla RAG \citep{rag-lewis} & \no & \no & \no & \no & \no & \no & \no \\
GraphRAG \citep{graphrag} & \no & \emph{communities} & \no & \no & \no & \no & \no \\
RAPTOR \citep{raptor} & \no & \emph{tree} & \no & \no & \no & \no & \no \\
LLM judge re-ranker \citep{rankgpt} & \no & \no & \no & \emph{ordering} & \no & \no & \no \\
Self-RAG \citep{selfrag} & \no & \no & \no & \emph{generator} & \emph{partial} & \no & \no \\
CRAG \citep{crag} & \no & \no & \no & \emph{corrective} & \no & \no & \no \\
DeepRead \citep{deepread} & \no & \emph{structure} & \no & \no & \no & \emph{agentic} & \textit{n/r} \\
ECM platforms \citep{paivarinta-ecm-2005,tyrvaainen-ecm-2006} & \textit{n/r} & \emph{lifecycle} & \textit{n/r} & \no & \no & \no & \textit{n/r} \\
Public KBs, DBpedia/YAGO \citep{auer-dbpedia-kb-2007,hoffart-yago2-kb-2013} & \textit{n/r} & \emph{ontology} & \no & \no & \no & \no & \no \\
\midrule
QDCVR (this work) & \yes & \yes & \yes & \yes & \yes & \yes & \yes \\
\bottomrule
\end{tabular}
\end{table}

Each tabulated system contributes one or two of these dimensions: GraphRAG
organizes a corpus as a graph, RAPTOR as a tree, Self-RAG and CRAG verify with
the generator, LLM judges order candidates, and ECM frameworks govern content
across its lifecycle. None combines a quality-gated acquisition pipeline, a
governed organization and lifecycle layer, first-class operational experience,
per-citation content verification, enforced abstention, budgeted multi-lane
retrieval, and a standard agent tool protocol.

This paper contributes that combination as one design: contracts that hold
across acquisition, organization, verification, and abstention
(Sections~\ref{sec:acquisition}--\ref{sec:abstention}), rather than capabilities
implemented per component. It does not propose a better ranker. On the
gold-annotated corpus BM25 reaches Hit@5 of 1.00 while the judged hybrid reaches
0.44; the value of the gate measured here is auditable, fail-closed,
per-citation provenance, not ranking quality (Sections~\ref{sec:e4} and
\ref{sec:inflation}).

\section{System overview and design principles}
\label{sec:system}

QDCVR (v2.3) is a self-hosted knowledge-base management system that turns a
folder of PDFs, Office documents, and scans into knowledge bases that agents
can be trusted to answer from. It is deployed as four cooperating services
behind one configuration file: a Python service that schedules parsing and
owns the OCR engine, a TypeScript web application that serves the interface
and acts as the server-side proxy, a Model Context Protocol (MCP) server that
exposes the platform's capabilities as tools \cite{mcp}, and a set of storage
engines that hold documents, vectors, graph edges and task state. The rest of
this section states what the system has to do, the principles that shaped it,
how it is put together, what it stores, and how its objects move through their
life cycle.

\subsection{Requirements}
\label{sec:req}

The requirements below were derived from the deployment rather than
postulated; each one names the failure it prevents.

\begin{description}
\item[R1 --- Parse heterogeneous enterprise documents.] Process regulations,
failure reports, scanned inspection records and journal articles arrive as
PDFs, Office documents and images. The system must convert them to text it
can govern, including scanned pages, tables and formulas, and it must make a
parse failure visible rather than storing an empty or garbled document. The
measured failure rate on a 100-paper acquisition was 21\% and is reported as
a first-class result (Section~\ref{sec:e1}).
\item[R2 --- Organize by content, not by folder.] A document's home must be
decided by what it says. Library membership, tags, descriptions and graph
edges are derived from the content, because folder paths encode who filed a
document rather than what it contains.
\item[R3 --- Manage the collection over its lifetime.] Libraries and
documents are created, moved, renamed, merged, split, deleted and re-indexed.
Every such operation must preserve agreement between the document on disk,
the tree registry and the library metadata, and the agreement must be
verifiable after the fact.
\item[R4 --- Retrieve with several independent channels.] Lexical, dense,
catalog and graph channels fail differently, so the system must be able to
generate candidates with each of them, under explicit budgets, and expose
which channel proposed what.
\item[R5 --- Do not promote a similarity score into a citation.] A candidate
becomes an answer only after its fetched content has been judged to contain
direct evidence for the question. When no candidate survives, the system
returns an explicit not-found report instead of the closest text.
\item[R6 --- Be operable and engine-agnostic.] The platform runs on a single
local workstation (RTX 5060, 8\,GB, Windows), degrades rather than crashes
when a model or a subprocess fails, exposes long-running work as pollable
tasks, and drives whichever agent harness the operator has configured rather
than being bound to one vendor.
\end{description}

\subsection{Design principles}
\label{sec:principles}

Six principles shape the architecture. They are stated as commitments because
each one is enforced somewhere in the implementation rather than left to
prompt discipline.

\begin{description}
\item[P1 --- Content-routed organization.] Documents are organized by what
they say, not by where they sat: placement, tagging, description and graph
edges are derived from the parsed content, and retrieval starts from a
library catalog rather than a directory path (\S\ref{sec:organization}).
\item[P2 --- Content outranks similarity.] No similarity score, however high,
promotes a document into an answer. A judging model must read the candidate's
fetched content and vote, and the platform fetches that content itself so an
agent cannot argue for a document it has not read (\S\ref{sec:gate}).
\item[P3 --- Every module terminates at a contract.] Ingestion ends in a
checked post-condition, retrieval in a budgeted return value, verification in
a rubric score, and management in an audit. A module whose result cannot be
checked is a module that cannot be trusted.
\item[P4 --- Fail closed.] Missing text, an engine error, a timeout or an
unparseable verdict rejects the candidate rather than passing it, and a
document that cannot be retrieved after indexing is a failed ingest rather
than a warning.
\item[P5 --- Leave an audit trail.] Every candidate, verdict, cut decision,
tool call and lifecycle operation is recorded, so that a claim about the
system can be re-examined from evidence instead of re-argued from memory.
\item[P6 --- Engine-agnostic agency.] Answer agents are pluggable harnesses
(14 adapters, including \texttt{omp}, Claude and ACP-compatible CLIs); the
platform constrains them through tool whitelists, budgets and prompts, never
by coupling to a particular vendor.
\end{description}

\subsection{Architecture and subsystems}
\label{sec:architecture}

Fig.~\ref{fig:arch} shows the four layers. The \emph{client and agent layer}
contains the web interface and the chat endpoint that adapts external agent
harnesses; requests are annotated with a lane policy before they leave this
layer. The \emph{tool layer} is an MCP server that exposes the platform's
capabilities as 95 tools in eight groups (knowledge-base CRUD and catalog,
document read/write, file system, parsing and task status, tags, search and
vector index, knowledge graph, and experience and persona management);
the lane policy is a server-side whitelist, so a lane physically cannot call
a tool outside its set. The \emph{domain service layer} holds the tree and
catalog services, the parsing scheduler that manages the OCR engine as a
supervised subprocess, the vector and two-stage candidate services, the
experience service, the graph service and the resident judging worker. The
\emph{storage layer} holds the document tree with its registries, the vector
collections, the graph database and the task and session databases.

Three properties of this arrangement matter for the rest of the paper. First,
the retrieval lanes of Section~\ref{sec:lanes} are not separate programs:
they are policies over one tool layer, which is why a new capability is
available to every lane at once and why lane behaviour is comparable across
harnesses. Second, long-running work (parsing, indexing, graph construction,
distillation) is exposed as a pollable task rather than a blocking call, so a
slow operation cannot stall an agent's turn. Third, the judging service is a
resident worker for engineering reasons, not architectural elegance: loading
the verifier per request cost 10--12\,s against 1--2\,s for a warm verdict,
which is why the worker, its file lock and its verdict cache exist
(\S\ref{sec:impl}).

\begin{figure}[t]
\centering
\includegraphics[width=\textwidth]{fig1_architecture}
\caption{QDCVR architecture. Documents enter through the acquisition
pipeline (left) and are parsed, quality-gated, split, described and indexed
into four stores. At run time (right), clients and external agent harnesses
reach the platform through the chat endpoint, which attaches a lane policy;
the policy is enforced against the MCP tool layer, whose 95 tools call domain
services that operate the storage engines. All retrieval lanes terminate at
the content-verification gate, and a lane that survives no candidate returns
an explicit not-found report. Organization and life-cycle management
(bottom) maintain agreement between the tree registry, the document files and
the library metadata.}
\label{fig:arch}
\end{figure}

Table~\ref{tab:modules} maps each function of the system to the components
that implement it, the design decision that shapes it, and the experiment
that provides evidence for it. The table is the paper's reading guide:
Sections~\ref{sec:acquisition}--\ref{sec:implementation} describe these rows
in order, and Section~\ref{sec:eval} evaluates them.

\begin{table}[!htbp]
\centering
\footnotesize
\setlength{\tabcolsep}{3.5pt}
\caption{System functions, their implementing components, the decision that
shapes each one, and the evidence that covers it. Evidence identifiers
E1--E7 refer to the evaluation subsections of Section~\ref{sec:eval}.}
\label{tab:modules}
\begin{tabular}{p{2.6cm}p{4.0cm}p{4.6cm}p{1.9cm}}
\toprule
Function & Components and representative tools & Key design decision & Evidence \\
\midrule
Acquisition and parsing & parsing scheduler, OCR subprocess, duplicate
detector, split gate & quality-gated: a document that cannot be parsed is
rejected visibly, and a document that cannot be retrieved afterwards is a
failed ingest & E1 \\
Knowledge organization & content router, catalog, tag registry, description
generator & content decides placement; the catalog is the first retrieval
instrument & E1, E2 \\
Life cycle and integrity & tree registry, library metadata, re-index tasks,
three-source audit & re-indexing is explicit; every operation is audited
afterwards & E2 \\
Document graph & graph projection of documents, libraries and tags;
related-document queries & the graph is a navigation index, never an evidence
source & E2 \\
Catalog-first navigation & library and document listings, progressive
disclosure & metadata before bodies, so a candidate set is chosen without
loading content & E4, E5 \\
Content-verified retrieval & vector, lexical, catalog and graph channels;
lane policies; retrieval gate & verification is a promotion condition, not a
ranking hint; verdicts fail closed & E4, E6 \\
Operational knowledge & experience store, credibility tiers, distillation
runs & a lesson is a claim about practice, so it is drafted, gated and
capped until it is vetted & E3 \\
Agent interface and operations & MCP tool layer, chat endpoint, lane
whitelists, task registry & every capability is a tool with a schema; long
work is pollable; engines are interchangeable & E5, E7 \\
\bottomrule
\end{tabular}
\end{table}

\subsection{Data model and storage layout}
\label{sec:datamodel}

The platform deliberately keeps its data model readable: documents are
Markdown files, and the two registries are JSON and YAML files that a human
can open. Table~\ref{tab:datamodel} lists the stores and the fields that the
rest of the paper refers to.

\begin{table}[!htbp]
\centering
\footnotesize
\setlength{\tabcolsep}{3.5pt}
\caption{Stores and their principal contents. Field names are the ones used
by the implementation and by the audit tools.}
\label{tab:datamodel}
\begin{tabular}{p{3.1cm}p{10.4cm}}
\toprule
Store & Contents \\
\midrule
Tree registry (\texttt{.tree-fs.json}) & version, root identifier and name,
timestamps, node counter and aggregate statistics; per node:
identifier, name, type, parent, path, children, metadata map, timestamps.
Authoritative structure of the collection. \\
Library metadata (\texttt{.knowledge-base.yml}) & library identifier, name,
description, creation and update timestamps, per-library configuration
(including distillation settings); per document: path, name, description,
tags, content hash, file size, vector-index metadata, free-form metadata.
The retrieval index read by catalog search. \\
Document files & Markdown body per document or per split part, plus an image
directory for figures extracted during parsing. \\
Vector store & one collection per library (cosine space, 1024-dimensional
embeddings, HNSW index), holding segment embeddings of 500 characters with
50-character overlap. \\
Graph store & nodes for documents, libraries, tags, folders and entities;
edges for membership, tagging, relatedness, identity and containment. \\
Task and session databases & pollable task records for parsing, indexing,
graph builds and distillation; chat sessions; service credentials;
distillation signals. \\
\bottomrule
\end{tabular}
\end{table}

Two consequences of this layout appear throughout the paper. Because the
registries are files rather than a database, a document's address is a path
and its identity is a path plus a content hash, which is what makes the
three-source audit of \S\ref{sec:org-registries} possible. And because the
vector store is a separate engine, it can disagree with the registries; the
system treats that disagreement as a detectable, repairable condition rather
than as an impossibility.

\subsection{Object life cycle}
\label{sec:lifecycle}

A document moves through five states. It is \emph{raw} when it arrives, with
a content hash recorded. It is \emph{parsed} when text has been extracted and
has passed the parse-quality checks. It is \emph{stored} when the body, its
description and its tags have been written to disk and registered in both
registries. It is \emph{indexed} when the vector store and the retrieval
indexes contain it, which requires an explicit indexing call. It is
\emph{verified} when a retrieval probe returns the document for a query
derived from its own content. Only the last state counts as a successful
ingest; a document that stalls at \emph{stored} is an incomplete operation,
not a warning.

Libraries have a simpler life cycle: a library is \emph{active} while it is
addressable and searchable, and it becomes \emph{merged} or \emph{split} when
an operator consolidates or divides it. Library metadata travels with the
library when it moves, including its distillation configuration, which is why
the configuration is stored per library rather than in the deployment file.
Platform-level versioning is explicit: the platform reports its version
(2.3.0 in this deployment) and can update itself from its repository through
the same command-line interface that performs the other operations.

Re-indexing deserves its own paragraph, because it is where the object life
cycle meets retrieval quality. Re-indexing is explicit, asynchronous and
observable: it returns a task identifier, it can be forced, and its
post-condition is a segment count that matches the document content. Nothing
in the write path re-indexes as a side effect. This is a deliberate trade of
convenience for the property that no document can silently be missing from
the vector store: the cost is that a writer who forgets the call leaves a
stale index, and the mitigation is the audit of \S\ref{sec:org-integrity},
which detects exactly that condition.

\section{Knowledge acquisition: parsing, ingestion, and index verification}
\label{sec:acquisition}

Acquisition is where a document becomes knowledge the system may later
answer from, and it is the stage this design treats most defensively.
The pipeline runs its gates in a fixed order, and a document that fails
one is refused rather than stored and indexed anyway. This section
follows the stages in execution order, with the audit numbers we have
and the measurement gaps we do not.

\subsection{Multi-format intake and parsing}

Upload starts in the browser at \texttt{POST /api/parse/file-vt}. The
Nuxt server forwards the payload to the backend parser endpoint
\texttt{POST /api/v1/parse/file/vt}, which is asynchronous and returns a
task identifier rather than a result. The backend then runs
MinerU~\citep{mineru} as a hidden subprocess on an ephemeral port to
convert the source into structured Markdown: OCR for scanned pages,
layout-aware recovery of tables, and formula extraction. The subprocess
is reachable only through the manager object, which resolves the
assigned port at runtime; the service never hardcodes it. Its combined
standard output and standard error are redirected to
\texttt{backend/logs/mineru-api.log} and never to a pipe---an earlier
build failed on pipe closure with a bare I/O error, and appending to a
file removed the failure mode. Parse jobs are polled with a 600\,s
timeout, so a stuck conversion terminates in a recorded failure instead
of blocking the caller.

The backend returns a \texttt{markdown\_path}---a filesystem location,
not document text. The proxy layer reads that file back and writes the
document into the knowledge base through the tree filesystem service,
which updates the on-disk Markdown, the tree index, and the per-library
metadata file in one operation. Keeping the transfer as a path rather
than a response body matters for large scanned documents: the parse
output can exceed what a single JSON response should carry, and the
proxy already needs filesystem access to place extracted images.

\subsection{Quality gates before storage}

Before any parse is scheduled, exact-duplicate detection hashes the
source with SHA-256 and compares it against the stored records.
Near-duplicates are then detected in vector space:
\texttt{kb\_find\_duplicates(threshold=0.90)}. A manual fallback channel
runs the same check at a 0.85 similarity threshold, combining a
filename and file-size match with a vector probe when the tool is
unavailable. On the 322-part corpus the audit recorded 0 true
duplicates and 35 near-duplicate pairs, all of them adjacent parts of
the same paper, which is the expected shape when a split document is
compared against itself.

The parse-quality gate inspects the first 1{,}500 characters of the
markdown and refuses the parse on any of five signatures: OCR garbage
(more than 30\,\% garbled), binary residue such as null bytes or base64
fragments, headings without body text, whitespace collapse, or a
language mismatch between the source and the extracted text. Refusal is
terminal for that attempt; the operator either re-parses with a
different OCR mode or reports the source as unusable.

Storage has one rule: parsed documents are written with
\texttt{kb\_doc\_save\_parsed}, never \texttt{kb\_doc\_create}. The
parsed path transfers the complete Markdown and copies all extracted
images; the generic create path truncates long content and drops
images, so using it on a parsed document silently loses both. Parts
produced by the split gate are stored individually through the create
path because each part is already a complete, size-bounded unit.

\begin{table}[t]
\centering
\caption{Acquisition gates, their trigger, and what refusal means. Gate
names follow the packaged ingestion contract.}
\label{tab:acq-gates}
\small
\begin{tabular}{@{}llp{5.5cm}@{}}
\toprule
\textbf{Gate} & \textbf{Trigger} & \textbf{On refusal} \\
\midrule
Exact duplicate & SHA-256 hit & Skip the source; report the existing path \\
Near duplicate & Vector similarity $\geq 0.90$; manual channel $\geq 0.85$ & Skip; supplement tags or description only \\
Parse quality & Garbled, empty, binary, or language-mismatched body & Reject the parse; re-parse or report the source \\
Split ceiling & Body beyond 30{,}000 characters & Split into parts; never store the oversized original \\
Atomic unit & Indivisible unit larger than the target & Keep whole with \texttt{oversized\_atomic\_unit}; hard cuts disabled \\
Tag anchoring & Tag absent from the read body sample & Return to extraction; do not release the document \\
Description readback & Claim not supported by the body & Rewrite the description; never the body \\
\midrule
Index verification & Collection wrong or \texttt{total\_chunks} $= 0$ & Ingest incomplete; do not report success \\
\bottomrule
\end{tabular}
\end{table}

\subsection{Structure-aware splitting and addressability}

Documents longer than a 30{,}000-character ceiling do not enter the
knowledge base as single files. A splitting pass proposes boundaries
from the document's own structure and a validator enforces a target
utilisation of 0.85 and refuses oversized atomic units: a block that
cannot be divided without cutting an equation or a table stays whole.
Each surviving part is written as its own file labelled
\texttt{part $k$ of $N$}, and every part is read through a line-offset
address rather than from its beginning, so retrieval can open a section
without paying for the whole document. The ceiling itself is a budget,
not a format limit. It keeps a stored unit inside the context window of
the agents that later read it, and it bounds the per-query read cost of
a single hit.

The measured outcome for the 100-paper arXiv corpus is 322 stored
parts, which is the corpus size the audits and the regression probes
later operate on. A separate 50-paper campaign under the same
30{,}000-character setting produced 165 documents (nine whole, 156
parts) and 3{,}798{,}446 characters from 2{,}672\,s of parsing; the
same campaign at a 2{,}000-character setting produced 2{,}713
fragments, which is what the ceiling exists to prevent. The vector
store chunks what is stored at 500 characters with 50 characters of
overlap~\citep{chroma,hnsw}; a separate dense replica index for
comparison uses 800-character chunks and produced 4{,}528 chunks for
the same corpus.

\subsection{Content-derived descriptions and tags}

Descriptions and tags are generated from the body text that the
ingestion agent has actually read, not from filenames or directory
positions. A description must carry four elements---subject,
method or technology, scenario or problem, and key data or
conclusion---with at least two concrete nouns among them, and every
claim is read back against the body before storage; a mismatch rewrites
the description, never the body. Tags pass a blacklist filter, are
normalised for case and language, and are then re-checked against a
sample of the body: a tag that does not appear in the text is dropped
rather than kept for vocabulary reuse.

This gate exists because of what an earlier audit found. Of 26 document
descriptions checked against their own bodies, 24 did not match: the
boilerplate ``Gutenberg front/back matter'' had been applied to real
chapters across a batch, and the librarian that reads part labels
selected two parts instead of seven. The wider catalog carries the same
class of defect, with 156 of 322 documents flagged for a broken part
prefix. Description quality is therefore an empirical result of this
work, not an assumption: a description is a claim about content, not a
fact about it. The consequence is architectural rather than editorial.
Because a claim can be wrong at ingest scale, it must be cheap to
verify and cheap to repair---which is why verification later became a
platform primitive (Section~\ref{sec:lanes}) and why the acquisition
audit ends with descriptions upgraded 165/165 rather than with a
generation count.

\subsection{Index construction and retrievability verification}

Indexing writes to three channels: a per-library vector collection, a
lexical BM25 index~\citep{bm25} over chunked documents, and a document
graph in Neo4j~\citep{neo4j}. Ingestion success is judged by observable
effects rather than by the tool's success envelope. A path that reported
success once returned an empty index payload after a GPU
out-of-memory failure, so the verification gate now requires a
non-empty collection name, the collection identifier expected for the
target library, and \texttt{total\_chunks} $\geq 1$ before the document
counts as indexed. Graph construction is asynchronous and returns a
task identifier that is polled, but task status is not accepted as
completion evidence---an earlier graph build left a task row stuck in
\texttt{running} while its artifacts were already queryable. The
step closes with an eight-point final check that includes the stored
copy read back at 500 characters, the chunk and graph node counts, and
three-layer metadata consistency across the tree index, the library
metadata file, and the files on disk.

\subsection{Asynchronous execution, failure recovery, and observability}

Parsing and graph construction are non-blocking: each returns a task
identifier immediately, and a task registry resolves that identifier
across tool processes, so a caller polls status instead of holding a
connection open. Batch ingestion is resumable, and it was resumed in
practice: an interruption during a 100-paper run was recovered by
restarting the client, which skipped documents already present and
continued from the recorded position. Per-paper parse time in the
50-paper campaign ranged from 16.3\,s to 40.4\,s over 26 subject
domains, 1{,}115.5\,s in total---an average of 22.3\,s per paper.

The same campaign exposes the limit of that bookkeeping. Of 100
submitted papers, the state file records 91 names, 79 were retrievable
in the knowledge base at evaluation time, and 21 never became
retrievable. The reported failure mode is formula-dense scanned
documents, and a mid-run client crash was recovered by restart, but no
artifact counts parse attempts, parse failures, or the error text of a
failed document. The 91-versus-79 difference is therefore an
observability gap: we can report what the artifacts show, and we cannot
attribute the difference from them.

Cross-lingual ingestion shows the same gates holding without
language-specific tuning. The corpus is 400 Wikipedia pages per
language in English, Chinese, and Japanese, ingested in two rounds; the
second round was bit-identical to the first
(Table~\ref{tab:acq-crosslang}). Ingestion success rate was 1.0 for all
three languages in both rounds. Membership asks where each page landed,
and the recorded knowledge-base membership accuracy is 0.875 for
English, 1.0 for Chinese, and 1.0 for Japanese. Storage completeness sums the
split parts of large documents back to the source length and is 0.975,
0.9559, and 0.9402 respectively. Self-retrieval Hit@1---the score for
the page's own indexed text---is 0.975, 1.0, and 1.0.

\begin{table}[t]
\centering
\caption{Cross-lingual ingestion, 400 Wikipedia pages per language,
first round. The second round reproduced these values bit-for-bit, so
one column is shown. Self-retrieval Hit@1 uses 40 queries per
language.}
\label{tab:acq-crosslang}
\small
\begin{tabular}{@{}lcccc@{}}
\toprule
\textbf{Language} & \textbf{Ingestion} & \textbf{KB membership} & \textbf{Storage} & \textbf{Self-retrieval} \\
 & \textbf{success} & \textbf{accuracy} & \textbf{completeness} & \textbf{Hit@1} \\
\midrule
English (en) & 1.0 & 0.875  & 0.975   & 0.975 \\
Chinese (zh) & 1.0 & 1.0    & 0.9559  & 1.0 \\
Japanese (ja) & 1.0 & 1.0    & 0.9402  & 1.0 \\
\bottomrule
\end{tabular}
\end{table}

Taken together, the gates buy a specific property: a document that
fails a gate is refused rather than silently indexed. The 100-paper
pipeline audit records that property directly---100/100 papers parsed
and routed by content, nine per-paper ingest checks 100/100, content
and storage audits 322/322 and 33/33. What the gates do not buy is
accounting. The artifacts carry no counter of attempted or failed
acquisitions, so the difference between 91 recorded and 79 retrievable
papers cannot be resolved into parse losses and lost records. We report
it as an observability gap, and we treat a submission as knowledge only
once verification has found it.

\section{Knowledge organization and lifecycle management}
\label{sec:organization}

Acquisition decides what enters the system; organization decides what the
system is. In this platform organization is not a folder convention that a
human maintains by hand, nor an emergent property of embedding similarity.
It is a governed process with three invariants: content decides where a
document belongs, the disk, the tree registry and the library metadata must
agree about every document, and every operation that removes or relocates
knowledge is preceded by a survey and followed by a verification. This
chapter describes how those invariants are implemented, which of them are
enforced by the interface rather than by convention, and where the current
implementation still allows the three stores to diverge.

\subsection{Content-routed placement}
\label{sec:org-placement}

Placement reads the parsed document, not its filename or its folder of
origin. For each arriving document the pipeline selects the library whose
stated scope covers the content, or proposes a new library when no existing
scope does; it then writes the document, derives a description and a tag
set, and runs the same final check used by acquisition before the document
counts as placed. Two rules keep the catalog usable as a retrieval
instrument. First, an existing library is preferred over a new one: the
catalog stays shallow, and the number of libraries an agent must inspect
before searching stays small. Second, a document is placed once; re-routing
is an explicit lifecycle operation, not a side effect of a later ingest.

The library tree supports one level of content grouping in which a parent
library contains child libraries, and the metadata of a child is addressable
through the parent. The implementation carries a caveat that we report
rather than hide: the vector segments of child documents are stored in the
parent library's collection, so a search that targets a child library
identifier returns no candidates. The shipped operational guidance is
therefore to search the parent collection and to use the child libraries for
navigation and scoping rather than as search endpoints. This is a limitation
of the current storage mapping, not of the routing rule, and it is the kind
of coupling between an organizational decision and a retrieval path that a
system paper has to expose.

\subsection{Catalog and metadata registries}
\label{sec:org-registries}

Two registries describe the same collection, and Table~\ref{tab:datamodel}
lists their fields: the global tree (\texttt{.tree-fs.yml}) is the
authoritative structure, and each library carries a
\texttt{.knowledge-base.yml} that indexes its documents. Documents themselves
are Markdown files with an accompanying image directory.

Agents read this catalog through a progressive-disclosure interface: the
library listing returns names, descriptions and document counts; the
document listing returns metadata without bodies; only an explicit read
returns content, and a read is attributed to a specific document. This design
choice is what makes the catalog-first retrieval lane viable: a caller can
narrow a search to a handful of candidate documents without ever loading a
body, which is why the catalog lane completes in 5--6 turns and 54--91~s
(Table~\ref{tab:lanes}).

Correctness of the catalog is defined as three-source consistency: for every
document, the file on disk, its entry in the tree registry, and its record in
the library metadata must agree on existence, path, size and index state. The
platform verifies this agreement rather than assuming it, and the audit is
productive: in the acceptance campaign a renamed document kept its old vector
segments, so the recorded chunk count read 18 where 9 chunks existed, and the
document matched its own former path at 0.9738 similarity, which is high
enough to be reported as a near-duplicate pair. Only an explicit force
re-index repaired it (18~$\rightarrow$~9). The design lesson is that metadata
consistency and index consistency are two different obligations: the first
can be satisfied by a single write operation, the second cannot, because the
vector store is updated by a separate call.

\subsection{Tag taxonomy and governance}
\label{sec:org-tags}

Tags are the coarse, human-legible axis of the collection. A tag registry
holds the vocabulary; each document's tags are normalized before they are
written, filtered against a blacklist of structural strings (section
headings, generic words such as ``introduction''), and re-checked against the
document body so that a tag must actually occur in the content it labels.
Retrieval by tag is a first-class operation, which makes tag quality a
retrieval concern rather than a cosmetic one.

Governance is a maintenance loop rather than a schema. The registry is
curated, not centralized: no taxonomy is imposed on authors, and the
vocabulary grows with the corpus. The audit therefore treats the registry as
a health surface: it reports tags with zero referencing documents, which
accumulate from test documents, abandoned naming experiments and chapter
headings, and it removes them only when instructed. In the audited campaign
165 documents passed the tag gate with no failures, and the live registry
held 316 tags. We do not claim a controlled vocabulary; we claim that the
vocabulary is inspectable and that orphaned entries are detectable.

\subsection{Duplicate detection and collection hygiene}
\label{sec:org-dedup}

Acquisition refuses exact duplicates by content hash and flags near
duplicates by vector similarity; at the collection level the question is
different, because the same content can legitimately arrive through two
ingest paths or into two libraries. The hygiene procedure is
survey-then-act: enumerate candidate duplicate groups at a similarity
threshold (0.90, with a manual fallback channel at 0.85), inspect each group,
and delete nothing automatically.

The audited run illustrates why the survey step matters more than the
threshold. On the 322-document corpus the detector reported 35 candidate
groups with similarities between 0.90 and 0.975; every one of them turned out
to be a pair of adjacent parts of the same long paper rather than a true
duplicate, and no deletion was performed. A policy that deleted on detection
would have destroyed 35 legitimate document parts. The same run confirmed
zero true duplicates, which is the expected outcome when acquisition gates
run before storage. Cross-library duplication is detected but not merged
automatically: merging two libraries is a lifecycle decision with provenance
consequences, so it stays with the operator.

\subsection{Library lifecycle operations}
\label{sec:org-lifecycle}

Table~\ref{tab:org-lifecycle} lists the operations the platform exposes, the
interface that performs them, and the consistency obligation each one
carries. All of them are available to agents through the same tool surface
that the user interface uses, which is deliberate: an agent that may read a
collection should be able to maintain it under the same rules a human
operator follows.

\begin{table}[t]
\centering
\small
\caption{Library and document lifecycle operations, their consistency
obligation, and the state observed in the acceptance campaign. ``Audited''
means the three stores were compared after the operation.}
\label{tab:org-lifecycle}
\begin{tabular}{@{}p{0.22\textwidth}p{0.24\textwidth}p{0.20\textwidth}p{0.24\textwidth}@{}}
\toprule
Operation & Interface & Obligation & Observed \\
\midrule
Create library & library create & tree + metadata + directory & audited \\
Rename or move document & document move & path updated in all three stores & audited; stale segments until force re-index \\
Update description or tags & document update & metadata rewritten, body re-checked & audited (165/165) \\
Update content & document update & new hash, re-index required & re-index not implicit \\
Delete document & document delete & removal from all three stores & audited (1~$\rightarrow$~0) \\
Batch delete & batch delete & per-document verification & 1/1 in campaign \\
Move library & library move & subtree and metadata move together & audited (0~$\rightarrow$~1) \\
Merge libraries & library merge & provenance preserved & demonstrated, not controlled \\
Split library & library split & no part orphaned & analysed, deliberately not executed \\
Rebuild index & re-index (task id) & segment count matches content & 18~$\rightarrow$~9 after force \\
\bottomrule
\end{tabular}
\end{table}

Two properties of this table carry the argument of the chapter. The first is
that re-indexing is not implicit. Creating, moving or rewriting a document
does not update the vector store; an explicit index operation does, and the
audit treats a missing or stale index as an incomplete operation rather than
as a warning. This is the single rule we would carry into any successor
system: an index that is updated by side effect is an index that silently
disagrees with the truth. Keeping an approximate index coherent with a
mutating document store is itself an active research problem
\citep{xu-spfresh-2023}; the contribution here is not a new maintenance
algorithm but a contract under which disagreement becomes detectable. The
second is the honest boundary of the evidence.
Move, rename, update, delete and re-index were exercised with real tool calls
and audited afterwards; merge and split are implemented and were demonstrated
on live libraries, but they have no controlled evaluation, and the
sub-library hierarchy was analysed and deliberately not executed pending an
operator decision. No concurrent organizational operation was tested. We
prefer to state this than to imply that a lifecycle operation is verified
because it has an interface.

\subsection{The document graph}
\label{sec:org-graph}

The platform maintains a property graph over the collection with document,
library, tag, folder and entity nodes and membership, tagging, relatedness,
identity and containment edges. Edges are derived from three signals:
co-membership of tags, vector similarity above a threshold, and the verdicts
produced by content verification during retrieval. Adding one document to a
library added four edges in the audited run (4{,}675~$\rightarrow$~4{,}679),
one per assigned tag; a full build over the 100-paper corpus produced 658
nodes covering 333 documents and 316 tags, and the 165-document intermediate
state held 335 nodes and 1{,}671 edges.

The graph has a deliberately restricted role. It is a navigation instrument:
it expands a candidate set, surfaces documents that share context across
libraries, and identifies central documents in a collection. It is never
evidence. Every candidate a graph traversal proposes still has to pass the
same content-verification gate as a candidate proposed by a dense retriever,
and the graph's own contribution to retrieval is measured separately: as a
standalone retriever it reached Hit@5 of 0.5625 on the industrial corpus and
0.2 on the second corpus, well below the lexical and dense channels
(Table~\ref{tab:retrieval}). We report this because it is the honest status
of graph-based organization in this system: useful for discovery, weak as a
ranker.

Two implementation gaps belong here as well. Graph construction is not yet
guaranteed by the ingestion path: a document ingested through one agent
harness received no graph node, so a library reported two of three
documents, which is a silent incompleteness of the kind the three-source
audit is designed to catch elsewhere. And batch graph builds have been
observed to leave a task record in the running state after the work
completed, so task status cannot be used as completion evidence; the graph
itself must be queried.

\subsection{Operational knowledge as a first-class collection}
\label{sec:org-experience}

Process documents describe what the organization intends; operational
lessons describe what actually happened. The platform stores the second kind
as a first-class collection with the same addressing, indexing and retrieval
interface as ordinary documents, plus two differences that reflect the
epistemic status of a lesson.

First, an experience carries a credibility tier, and the tier controls
whether it may be shown (Table~\ref{tab:exp-tiers}). Knowledge that is a
claim about practice is held to a stricter bar than a document: the
experience channel's similarity threshold is 0.55 against the document
channel's 0.35, and an unvetted or disputed entry is capped even when it
scores well. Second, entries are drafted before they are served. Extraction
produces candidates from ingested content or from recurring question-answer
pairs; automated distillation runs as a separate agent process on a schedule,
guarded by a circuit breaker (three consecutive failures suspend it for 24
hours) and a heuristic fallback whose output is marked low-confidence.
Nothing reaches the served collection without an explicit approval step,
which is the same fail-closed posture the retrieval gate applies to
citations.

\begin{table}[t]
\centering
\small
\caption{Experience credibility tiers and modifiers. Tiers combine retrieval
score, content-verification score, ratings and review count; modifiers cap
the tier rather than promote it.}
\label{tab:exp-tiers}
\begin{tabular}{@{}llp{0.46\textwidth}@{}}
\toprule
Tier & Condition & Effect \\
\midrule
P0 strong & sim $\ge 0.65$ $\wedge$ content $\ge 6$ $\wedge$ rating $\ge 4$ $\wedge$ reviews $\ge 1$ & may be presented first \\
P1 reference & sim $\ge 0.45$ $\wedge$ content $\ge 4$ & presented with credibility noted \\
P2 weak & sim $\ge 0.35$ $\wedge$ content $\ge 3$ & suppressed unless expanded \\
discard & content verification fails or sim $< 0.35$ & never presented \\
\addlinespace
disputed & reviews $\ge 3$ $\wedge$ rating $< 2.0$ & capped at P2 \\
unvetted & 0 reviews $\wedge$ 0 applications & capped at P1 \\
auto-extracted & not human-reviewed & capped at P1 until applied twice and rated $\ge 4$ \\
\bottomrule
\end{tabular}
\end{table}

The evidence for this subsystem is deliberately presented as case evidence.
One extraction campaign produced thirteen candidate drafts; one end-to-end
scenario sedimented a single experience, which a later global experience
search returned first with a local similarity of 0.711; and the approval path
was never exercised, because the draft pool was empty when it was tested. A
single scenario and thirteen candidates cannot support an effectiveness
claim, and we do not make one. What the campaign does establish is that
operational knowledge travels through the same pipeline as documents ---
addressed, indexed, gated and auditable --- rather than living in a wiki or a
prompt.

\subsection{Integrity verification as routine}
\label{sec:org-integrity}

Every property claimed in this chapter is checked by a verification pass that
runs against the live collection: three-source consistency per document,
document integrity, parse quality, index coverage with optional repair, tag
health, experience health, and a scorecard that summarizes the findings. The
pass is read-only by default; repairs require an explicit instruction, and
the report is produced before any repair. This is a deliberate ordering: an
audit that repairs as it goes cannot report what it found.

The pass is not ceremonial. It is how the two index-integrity incidents of
this deployment were found --- a silently lost collection and a stale index
--- and how the reorganization campaign was certified after the fact
(322/322 metadata checks, 165/165 description checks, 33/33 sampled
consistency checks, 15/15 retrieval regression probes with no degradation).
In the same campaign the audit also surfaced the failures reported above: the
rename-time chunk duplication, the missing graph node, and a task-status
record that outlived its task. A collection whose health is a command is a
collection whose health can be argued about with evidence.

\subsection{What organization buys, and what it does not}
\label{sec:org-summary}

Organization buys candidate quality and auditability. Because the catalog is
curated and consistent, an agent can select a library and a document set
without loading bodies, which is the mechanism behind the catalog-first lane
and behind cross-library discovery in an environment with nineteen libraries.
Because every placement, tag, description and deletion left a record, the
collection can be re-verified after the fact; the numbers in this chapter are
not survey estimates but audit outputs.

Three things organization does not do, and the paper does not claim. It does
not repair a bad parse: a document whose parse lost its tables cannot be
recovered by placing it well. It does not prevent divergence by
construction: the write path can leave the vector index behind, and only an
audit and an explicit re-index restore agreement. And it does not scale
without limits: the largest collection we report is 322 documents in five
libraries, the merge and split operations are unevaluated, and no
organizational operation was tested under concurrency. These are engineering
boundaries, and the next chapter describes the retrieval contract that is
built on top of them.

\section{Content-verified retrieval}
\label{sec:retrieval}
\label{sec:lanes}

Every citation a lane proposes must pass a content-verification gate before it
may appear in an answer. The gate scores each retrieved segment on a 0--8 rubric
(topic 0--3, scenario 0--3, evidence 0--2) and keeps a document only if at least
one of its segments reaches the acceptance threshold: 0.5 in general, 0.3 for
enumeration questions. It processes at most six references per call, takes the
matched text from the server side rather than from the caller, and fails closed.
No surviving evidence means the caller must answer that the corpus does not
contain the answer.

\subsection{One query through the gate}
\label{sec:trace}

Q2 of Section~\ref{sec:eval-setup} is \emph{``EL inspection shows edge delamination:
root cause and corrective actions''}, asked through the public chat API with the
failure-analysis library pinned and the Hybrid lane selected. Its timeline shows
where candidates come from, what the gate decides, and how much text the agent
may read.

\begin{enumerate}
\item \textbf{Policy.} The chat endpoint attaches the Hybrid lane's tool
whitelist and instructions server-side, so the agent cannot call a tool outside
it.
\item \textbf{Candidate generation (one call).} The agent issues a single
\texttt{kb\_hybrid\_search}. Two channels run server-side: a dense vector channel
(top-30, cosine threshold 0.35) and a catalog keyword channel over the pinned
library's descriptions. Candidates merge on (library, document), documents
agreed by both channels are ordered first, and the pool is capped at twelve. On
Q2 the pool is the gold root-cause report plus eleven same-domain siblings.
\item \textbf{The gate.} The server fetches each candidate's full text itself,
segments it, and scores \emph{does this text contain direct evidence that answers
the query} in batches of at most six. The agent never supplies text, so it cannot
argue for a document it has not read. On Q2 the gold report scores 0.846 while
five extruder-telemetry reports score 0.887--0.926.
\item \textbf{Cut and read budget.} A document-level relative cut keeps scores
$\geq \mathit{best}-0.15$ with a top-5 floor and a cross-domain anchor guard, so
an inflated out-of-domain score cannot set the cut line; the top five survivors
are returned in full, at most 2\,000 characters each. That read window is the
only full-content channel the agent sees; the rest of the pool survives as
scores in the audit trail.
\item \textbf{The answer.} Citations name the document path, the lane, and the
engine score, graded P0 (direct evidence) through P2 (weak).
\item \textbf{The terminal paths.} If the cut leaves no survivor, Lane~A
escalates once to Lane~B; if Lane~B also exhausts its budget, the run ends in an
explicit not-found report instead of a speculation.
\end{enumerate}

The trace also shows where the design fails. Under the shipped re-ranking role,
Q2's gold (0.846) is pushed to rank eight by the five higher-scoring telemetry
reports, outside the five-document read window, so the agent reads confusers
(Fig.~\ref{fig:anatomy}). Under the filter role the same pool, the same scores,
and the same verdicts leave the gold first: same gate, same evidence, only the
role of the induced ordering changes. This one trace foreshadows the paper's
central result (Section~\ref{sec:inflation}).

\subsection{The gate contract}
\label{sec:gate}

The gate is a locally hosted evidence-classification model, the Laya judging
service \cite{laya}, invoked through a \texttt{kb\_laya\_judge} tool. Given a
query and a batch of document references, the server fetches each document
itself, segments it, and scores every segment under one of three criteria:
\emph{evidence} (specific substantiation; a mere topic match is insufficient),
\emph{instance} (the text contains at least one member of the requested class,
routed automatically for enumeration questions), or \emph{auto}. A document
survives only if its best segment reaches the threshold, so a document at vector
score 0.40 with one \emph{yes} verdict survives while a document at 0.95 that the
judge rejects is dropped.

At the protocol layer the platform documents the score as a 0--8 rubric
(topic 0--3, scenario 0--3, evidence 0--2; $\geq$6 answers directly, 5 escalates to
the librarian lane, $\leq$4 drops). The $[0,1]$ score is the machine-readable
projection of that rubric. The projection can saturate --- the logged maximum
score is 1.0000 --- which is one reason thresholds carry little discriminating
power between golds and same-domain confusers (Section~\ref{sec:inflation}).

Every operating constant is a shipped production default: the thresholds, the
six-reference judge batch, the hybrid candidate cap of twelve, and the 0.15
relative-cut margin with its top-5 floor. They were fixed during earlier
deployments on a different corpus and were not tuned on the evaluation questions
of Section~\ref{sec:eval-setup}. The classifier is a single, locally fine-tuned
evidence classifier whose checkpoint reference ships with the platform
artifacts; it is deterministic in our replays (repeat calls with identical
inputs reproduced identical orderings), and every verifier finding in this paper
inherits that choice, a scope limit noted in Section~\ref{sec:threats}.

Fail-closed semantics apply to every failure mode: fetch errors, engine errors,
and timeouts count as rejection. Survivors are aggregated into an evidence pack
with per-passage provenance.

\subsection{Three lanes}

The three lanes are not three implementations of one idea; they are three points
on a latency/coverage frontier (Table~\ref{tab:lanes}; Fig.~\ref{fig:lanes}). The
Hybrid lane is the fastest and the most auditable, because it moves orchestration
server-side and ships its audit trail in one response; it is the interactive
default. The Librarian lane is the one that must not miss a document: it spends
its twelve-call budget on catalog navigation and buys exhaustive coverage with
explicit blind-spot disclosure. The Search lane is the latency-tuned interactive
path for pinned libraries.

\begin{table}[!htbp]
\centering
\scriptsize
\setlength{\tabcolsep}{3.5pt}
\caption{The three lanes as operating points. Measured figures are the
round-1 agentic runs of Section~\ref{sec:agentic}. All three share one
verification gate and one honesty contract; they differ only in retrieval
affordance and budget.}
\label{tab:lanes}
\begin{tabular}{p{1.85cm}>{\raggedright\arraybackslash}p{3.05cm}>{\raggedright\arraybackslash}p{2.15cm}>{\raggedright\arraybackslash}p{2.25cm}>{\raggedright\arraybackslash}p{3.15cm}}
\toprule
Lane & Optimises for & Tool budget & Measured (round 1) & Use when \\
\midrule
A. Search & latency on pinned libraries & 20-turn cap & 6--8 turns; 70--165\,s & interactive default; single-library questions \\
B. Librarian & exhaustive recall across all libraries & hard 12-call budget & 5--6 turns; 54--91\,s & must-not-miss; enumeration questions; no library pinned \\
C. Hybrid & single-call server-side orchestration & exactly 1 call & 3--5 turns; 36--70\,s & fastest grounded answer; latency-critical \\
\bottomrule
\end{tabular}
\end{table}

\begin{figure}[t]
\centering
\includegraphics[width=\textwidth]{fig2_lanes}
\caption{The three retrieval lanes. Every lane ends at the same
content-verification gate and the same honesty contract: zero verified
evidence escalates (Lane A$\to$B) or terminates in an explicit not-found
report.}
\label{fig:lanes}
\end{figure}

\textbf{Lane A (Search).} Phase A0 rewrites the query into (subject, attribute,
constraints) form, splits multi-entity and comparison questions into parallel
sub-queries, and routes incident-style queries to the experience store
(Section~\ref{sec:system}) first. Phase A1 issues one wide dense query (top-30,
threshold 0.35, cross-library balancing), deduplicates by document, and keeps all
candidate documents rather than truncating to the top chunks. Phase A2 submits
citation-only references (library id plus document path) to the gate in batches
of at most six, and Phase A3 returns the judge's evidence pack, the only
full-content channel in this lane, from which the agent composes an answer whose
citations carry confidence levels (P0 direct evidence, P1 partial, P2 weak).

\textbf{Lane B (Librarian).} A fixed spine under a hard budget of twelve tool
calls. L0 reads the library catalog (ids, names, descriptions, counts) and prunes
shelves by the rewritten query; L1 lists each surviving shelf's documents with
descriptions; L2 routes by query shape, sending literal and enumeration questions
(``which documents mention X'') to a machine-side scanner ($\approx$0.5\,s per
document, no LLM turns) and semantic questions to judging; L3 judges in batches
of at most six, because larger batches historically triggered transport timeouts
and fail-closed batch loss; L4 reads only survivors, bounded to 2\,000 characters
and starting at the judge-reported evidence line; L5 emits a five-part report
(search paths, answer, sources, confidence, blind spots). Two rules are absolute:
catalog descriptions never prune documents by themselves, and an exhausted budget
terminates in blind-spot disclosure, not in a guessed answer.

\textbf{Lane C (Hybrid).} Agentic latency is dominated by LLM turns, not by tools
(Section~\ref{sec:composition}), so Lane C spends exactly one call by moving the
whole pipeline into the server. A single \texttt{kb\_hybrid\_search} call executes
P1 (a parallel block of a vector channel, top-30 at threshold 0.35, and a catalog
keyword channel); P3 merges candidates by (library, document), agreed-by-both
documents first, capped at twelve; P4 judges every candidate with the gate in
batches of six (criterion \emph{evidence}, threshold 0.5); P4b applies the
document-level relative cut; and P5 reads the top five documents in full
($\leq$2\,000 characters each). The response carries the complete audit trail
(per-lane candidates, merge counts, judge statistics, cut decisions, read
contents), so a one-call lane is still an inspectable lane.

\subsection{Selecting a lane is a budget decision}

The three lanes trade latency for recall under an enforced call and turn budget,
so the choice is not which retriever is better but how much of the budget to
spend before the gate runs. Hybrid is the interactive default, Librarian the
must-not-miss lane, Search the pinned-library path.

What makes the lanes interchangeable is the gate. All three terminate in the same
contract --- one gate, one fail-closed semantics, one evidence pack --- so
replacing a lane changes the retrieval affordance and the budget, not the meaning
of a citation. Escalation follows a trigger rather than a hunch: zero verified
evidence escalates Lane~A to Lane~B exactly once. Query-shape auto-selection is
future work (Section~\ref{sec:threats}).

One caveat is honest to record. On the evaluated corpus Lane~A occupies no
\emph{distinct} operating point of its own: the Hybrid lane is faster and the
Librarian lane more exhaustive. A two-lane deployment (Hybrid plus Librarian) is
therefore defensible today, with the Search lane retained as the pinned-library
interactive default; because the lanes share one gate and one abstention
semantics, that reduction costs coverage, not trust.

\subsection{The gate as a platform primitive}

The gate is not a prompt. It is a server-side tool reachable only through the
platform's own interface, and its operating point is an explicit argument rather
than a sentence in an instruction:
\texttt{kb\_laya\_judge(query,\allowbreak documents,\allowbreak
criterion=auto,\allowbreak threshold=0.5,\allowbreak
max\_evidence\_chars=40000)}. The platform, not the agent, chooses the criterion
and fixes the threshold. Calls are serialized through a single judge worker
--- an early concurrent workload crashed the worker with a torch paging-file
error, so verdicts take a lock --- with a 900\,s call timeout, and every verdict
is logged with its query, references, scores, and decision.

That placement makes verification a property of the layer rather than an
aspiration of the agent. Any lane, harness, or external client that answers
through the platform inherits the gate; none can skip it, and none can
re-threshold it per query. Enforcement has a second half in the lane policy,
which attaches a per-lane tool whitelist server-side: a librarian-lane agent is
denied the vector and two-stage tools at the protocol level, so it cannot reach
around the catalog-first path it was given.

Externalized verification also distinguishes this design from verifiers that run
inside the generator. Self-RAG and CRAG verify during generation with the
generator itself \cite{selfrag,crag}; the gate sees only (query, document text)
pairs and is consumed as a tool. That is what makes it reusable and inspectable,
and, by the same placement, fallible: the LLM-judge critique applies here as much
as anywhere \cite{zheng-llmjudge}, which is why verdicts fail closed, are logged,
and are reported as evidence rather than as truth (Section~\ref{sec:inflation}).
The cost is paid once per deployment: a twelve-candidate verification adds
$\approx$0.5\,s to a call (Table~\ref{tab:retrieval}), tens of milliseconds per
verdict once the worker is resident, while a cold judge load costs 10--12\,s
(Section~\ref{sec:impl}).

\subsection{Honest abstention}

A threshold that never rejects a candidate cannot refuse a query. On four
out-of-corpus probes the judged hybrid still keeps 9--11 of its 12 candidates
(Section~\ref{sec:abstention}, Fig.~\ref{fig:ood}), and the graded probes behave
no differently: prompts with partial coverage and prompts with none both keep
10--12 of 12 candidates, with top scores up to 0.983
(Table~\ref{tab:robustness} and Section~\ref{sec:e7}). Because the $[0,1]$ projection saturates, no cut
point separates a supported answer from an unsupported one.

Abstention therefore lives in the reading-and-answering layer, enforced by the
lane contract rather than advised in a prompt. A caller with no surviving
evidence must report a not-found with its search scope: which libraries and how
many documents were searched, how the query was routed, and which topics sit
outside the corpus. The escalation rule makes that state terminal, so an
exhausted budget produces a report and not a guess.

Because abstention is a terminal state enforced server-side, it can be evaluated
like any other outcome: all six out-of-corpus probe runs returned explicit
not-found reports with zero fabricated citations, even though the tool layer had
handed them candidates (Sections~\ref{sec:agentic} and \ref{sec:abstention}). The
division of labour is the point: the gate decides what may be cited, the reading
layer decides whether to answer at all, and the measurements put the reliability
in the second and not the first.

\section{Agent interface, implementation, and non-functional properties}
\label{sec:implementation}
\label{sec:impl}

\subsection{The tool surface}
\label{sec:impl-tools}

The platform is exposed to agents as a single Model Context Protocol
server~\cite{mcp} carrying 95 tools in eight groups: knowledge-base CRUD and
catalog; document read/write; file system; parsing and task status; tags;
search and vector index; knowledge graph; and experience and persona
management. The count is the number of registered tool decorators in
\texttt{kb-mcp/server.py}, checked one-to-one against the function
definitions they register; the project's own documentation had drifted to a
lower number, so we report the measured one. The interface is not a reduced
agent-only view: it owns no private data path. Writes go through the same
HTTP services the user interface calls, reads read the same
\texttt{.tree-fs.json} and \texttt{.knowledge-base.yml} registries the web
server reads, and the MCP process itself contains no HTTP code, only a
client module.

Three governance rules make a surface this wide safe to hand to a model.
First, read-path tools and write-path tools are disjoint families, and
retrieval lanes are granted the read family only. Second, a stored parsed
document cannot enter a knowledge base through the generic document writer:
it must pass the parsed-document writer (\texttt{kb\_doc\_save\_parsed}),
the single entry point that carries parse provenance and forces the
ingestion quality gates to run. Third, long-running operations---a parse
batch, a re-index, a graph build, a meditation pass---return a task id
immediately instead of blocking the session, and are polled through the
task-status tools backed by an in-process task registry.

\subsection{Harness integration and lane policy}
\label{sec:impl-harness}

One chat endpoint adapts fourteen agent harnesses to a uniform message
stream, so a vendor CLI, an ACP-compatible agent, and the platform's own
engine all reach the same retrieval lanes. The lane policy of
Section~\ref{sec:lanes} is enforced here rather than by prompt discipline:
each lane is attached to a tool allow-list, of which 21 entries are
read-only MCP retrieval tools, and to a deny-list of more than 40
harness-internal tools (file editors, web search, sub-agent spawning). The
gate was observed working in the harness campaign: on its first turn a model
attempted to end the session early, the attempt was blocked by the
deny-list, and the retrieval that followed was unaffected. Sessions are
budgeted as well, with per-lane turn and call caps
(Table~\ref{tab:lanes}); the campaign ran 18 end-to-end scenarios, all of
which passed with no fabricated answer.

Three operational findings from that campaign matter to callers. A trivial
ping prompt triggered a six-skill pre-flight before any retrieval, and the
probe was killed by a 30\,s timeout; the platform's default caller timeout
is therefore 420\,s, not the few seconds a request such as ``list my
libraries'' suggests. The streaming endpoint emits bare \texttt{data:}
frames with no event names, so a client must parse data-only frames instead
of waiting for a named event. Finally, engine choice dominates cost by one
to two orders of magnitude at the same quality tier: \$0.90--1.80 per
question for the agent harness, \$0.11--0.36 for the second engine, and
\$0.03--0.07 for the third. Wall clock was 50--176\,s, 23--52\,s, and
34--171\,s respectively, but the third engine's engine-side time was only
15--83\,s: its extra wall clock is harness cold start of 50--90\,s per call,
not retrieval. One residual gap appeared in the same campaign: a document
ingested through the third harness had no graph node, while the same
operation through the agent surface created one.

\subsection{Implementation and runtime engineering}
\label{sec:impl-runtime}

The backend is Python 3.12 with FastAPI. ChromaDB~\cite{chroma} stores one
HNSW~\cite{hnsw} collection per knowledge base in cosine space; embeddings
are BGE-M3~\cite{bge-m3}, 1024-dimensional and normalized, encoded in
batches of 32 by a lazy-loaded singleton with an LRU encode cache. Neo4j
5.20~\cite{neo4j} holds the document graph, jieba~\cite{jieba} tokenizes the
keyword lane, MinerU~\cite{mineru} performs parsing and OCR, and MCP is
spoken over stdio and SSE.

The enforcement described in Section~\ref{sec:lanes} depends on engineering
invisible in the interfaces. The judging model runs as a resident worker
behind a machine-level file lock, with an LRU verdict cache and explicit
CUDA-cache release between workloads. A cold judge load costs 10--12\,s per
batch; the resident worker answers in 1--2\,s. The worker also removes a
failure mode met in production: when judge and bulk-embedding workloads ran
interleaved without cache cleanup, one end-to-end run degraded from 68\,s to
27 minutes under the Windows display driver model's paging behaviour. Judge
calls are serialized by one lock, because two concurrent judgments on a
single GPU raise a native allocator failure; a call is held to a 900\,s
ceiling and fails closed beyond it. MinerU
is launched as a hidden subprocess with platform-specific parent-death
cleanup---a Windows Job Object with kill-on-close, Linux
\texttt{prctl(PR\_SET\_PDEATHSIG)}, and a process group plus
\texttt{atexit} on macOS---and its stdout is redirected to a log file rather
than a pipe, because the pipe variant failed with a spurious
\texttt{Errno 22} once the reader closed. Before start, the backend probes
its own port with a bare socket bind and refuses to start if the port is
occupied, so a surviving process cannot be shadowed by a second one. Vector
writes are serialized per collection, after one production incident traced
to concurrent segment handles; the keyword index updates incrementally.
Every HTTP client is constructed with \texttt{trust\_env=False}, so a
corporate proxy cannot intercept calls to localhost.

\subsection{Configuration, deployment, and reproducibility}
\label{sec:impl-deploy}

A single \texttt{config.yml} is the source of truth for ports and
thresholds; environment variables override it. The mechanism is used in
practice, not only documented: in the harness campaign the services ran on
web 6789, backend 8771, and kb-mcp 8000, the backend port being an override
of the shipped development default. Rate limiting (a general bucket and a
smaller bucket for expensive routes), authentication tokens for the MCP,
user, and legacy clients, and the retrieval thresholds all live in the same
file. The \texttt{ragctl} command-line tool covers setup, start, status,
dependency, configuration, and update operations. The system runs
self-hosted on a single workstation: ChromaDB runs in-process, Neo4j runs
locally on bolt port 7687, and MinerU runs as a local subprocess, so no
external service is required for a deployment. The same file carries the
higher retrieval threshold (0.55) used by the experience subsystem, whose
rolling summary of operational lessons is regenerated by a rules-based pass
every fifth ingestion. Source code and all experiment scripts are public at
\url{https://github.com/kingdol666/rag-knowledge}.

\section{Evaluation}
\label{sec:eval}

We evaluate the deployed system on three corpora that differ in origin, language and scale, using three metric families, and we organise the evidence into seven groups. E1 and E2, reported here, cover the stages that precede a question: acquisition and parsing, and organisation, consistency and management operations. E3--E7, reported in Sections~\ref{sec:e3}--\ref{sec:e7}, cover retrieval-quality ablations, agent behaviour and verifier behaviour. Numbers are reported against the corpus that produced them, and results from different corpora are never merged.

\subsection{Experimental setup}
\label{sec:eval-setup}

\textbf{Corpora.} (i)~\emph{Deployed industrial corpus}: five knowledge bases in a live deployment hold 42 documents, one of them containing twenty near-duplicate equipment-telemetry reports that share prefixes, vocabulary and structure, a deliberate stressor for near-duplicate detection and confuser density~\cite{manku-near-dup-2007}. Sixteen gold questions and four out-of-corpus probes were run against it. (ii)~\emph{arXiv corpus}: 100 submitted papers spanning about 25 domains; at evaluation time the state file records 91 paper names and 79 papers are retrievable. Content routing placed them in five domain libraries holding 322 stored parts, and ten gold questions (BQ01--BQ10) target books referenced by these papers, keyed by arXiv identifier. The live environment held 19 libraries. (iii)~\emph{Cross-lingual ingestion corpus}: 400 Wikipedia pages per language for English, Chinese and Japanese, with 40 extractive cloze queries per language whose answers occur verbatim in the source page. It is used for ingestion only; we report no retrieval-quality numbers for it.

\textbf{Questions and gold annotations.} The sixteen industrial questions (Table~\ref{tab:questions}) were written by the authors rather than sampled from logs, a circularity we state rather than hide; each identifies one or two gold documents by a unique filename fragment and one to three fact markers. Q16 names the event \texttt{heatdecay092}, a token that also occurs in its gold document's filename, so a retriever ranking on filenames alone could match it without reading content; we keep the question and report the weakness. On the arXiv corpus a question passes only if the gold identifier is traceable in the answer or the tool trajectory and at least 60\% of the gold keywords occur at word boundaries, evaluated deterministically with no language model in the loop.

\begin{table}[!htbp]
\centering
\scriptsize
\setlength{\tabcolsep}{3pt}
\caption{Gold-annotated question set over the five production knowledge
bases. ``Gold'' lists unique filename fragments of the ground-truth
documents. All sixteen questions serve the tool-layer experiment; Q2, Q4, Q6, Q7 also
serve the agentic experiment.}
\label{tab:questions}
\begin{tabular}{lp{5.3cm}p{4.0cm}p{2.8cm}}
\toprule
\# & Question (abridged, translated) & Gold document(s) & Design intent \\
\midrule
Q1 & EVA encapsulation film lamination: how to set the temperature, pressure
and time window; how to control cure degree & EVA process essentials; EVA
hydro-thermal aging mechanism & multi-gold \\
Q2 & EL inspection shows edge delamination: root cause and corrective
actions & EL-2025-0317 delamination RCA report & incident $\to$ RCA \\
Q3 & Why does humid aging of EVA generate acetic acid and corrode ribbons;
moisture transport mechanism & EVA hydro-thermal aging mechanism & mechanism \\
Q4 & Moisture-barrier behavior of POE encapsulant; lamination differences
vs.\ EVA & POE process essentials & cross-material \\
Q5 & Summer power drop at a PV plant: systematic troubleshooting order &
Power-drop troubleshooting manual; O\&M lessons & synonym gap
(``drop'' vs.\ ``plunge'') \\
Q6 & Trigger condition, nominal value and tolerance of the line-1027 part
weight alarm & Watch-clause document (31.35\,g $\pm$0.05\,g, 5\,min) &
numeric exactness \\
Q7 & Melt temperature fell 206$\to$180$^\circ$C on extruder line 1: final
root-cause verdict & TC0 dead-loop RCA revision (among 20 telemetry dups) &
confuser density \\
Q8 & PET biaxial stretching process flow and critical temperature ranges &
PET process regulation & plain factual \\
Q9 & Frequent inverter alarm trips: handling steps & Inverter alarm handling
manual & PV O\&M family \\
Q10 & Hot-spot acceptance criteria and disposal & Hot-spot detection and
disposal procedure & PV O\&M family \\
Q11 & String mismatch causes; IV-curve diagnosis & String mismatch analysis
method & PV O\&M family \\
Q12 & PP cast-film extrusion parameters; chill-roll temperature & PP cast
film process regulation & plain factual \\
Q13 & Weight closed-loop run1: staircase excitation design & Weight
optimization staircase record & telemetry cluster \\
Q14 & FT full-function regression after P0 optimization & FT regression test
record & telemetry cluster \\
Q15 & Multimodal diagnosis: fusing scalar, vector and image signals &
Three-modality closed-loop record & telemetry cluster \\
Q16 & Post-compensation recovery verification of the heatdecay092 event &
heatdecay092 compensation verification (among 3 near-dup incident docs) &
confuser density \\
\bottomrule
\end{tabular}
\end{table}

\textbf{Metric families.} Retrieval results use Hit@5, Recall@5 and MRR@5, following passage-ranking evaluation conventions~\cite{craswell-msmarco-2021}. A graded measure, nDCG@5~\cite{jarvelin-cumulated-gain-2002}, was computed only inside one exploratory arm of the cross-lingual benchmark that we do not report; no nDCG@5 value appears here. Grounding and key-fact verification use deterministic string tests: whether the answer cites the gold document, whether at least 60\% of gold key facts appear at word boundaries, and whether part, section and table citation markers are present. These checks are mechanical, not model-based faithfulness estimates such as RAGAS~\cite{es-ragas-2024} or ARES~\cite{saad-falcon-ares-2024}. The third family covers wall-clock time, agent turns, tool calls, monetary cost, and the correctness of refused answers.

\textbf{Controls.} Deterministic lanes were re-run and compared field by field with latency excluded, so a re-run could disagree only on content and ordering; model-driven lanes were repeated with the same model version and are reported as mean with the observed spread. A fixed retry policy applied: an empty zero-turn response or a transient gateway timeout was retried once, which fired on two runs and succeeded both times. Before measurement we repaired two corpus-integrity faults through the platform's own operations: one document missing from its vector collection was re-indexed, and one library whose collection had been silently lost was rebuilt. All measurements ran on one workstation whose RTX~5060 (8\,GB) serves the embedding model and the judge worker, with every service local under Windows.

\textbf{What is deterministic.} The tool-layer retrieval lanes, the grading rule, the audits, duplicate adjudication and the cross-lingual ingestion pipeline are deterministic: repeating them reproduces the same values, and cross-lingual round~2 was bit-identical to round~1. The agent lanes, and anything derived from generated text, are not, and are reported with their repetition spread.

\subsection{E1: Acquisition and parsing}
\label{sec:e1}

One hundred submitted papers were parsed and routed by content into five domain libraries, producing 322 stored parts after the 30{,}000-character split. Table~\ref{tab:acq-pipeline} lists the acquisition audit: nine per-paper ingest checks passed 100/100, no true duplicate was found against 35 near-duplicate pairs, 165 document descriptions were upgraded, the content and storage audits passed 322/322 and 33/33, all 15 retrieval regression probes passed, and the dense replica index holds 4{,}528 chunks. Values are asserted by the run scripts from saved artifacts rather than hand-counted; the content audit samples the first 1{,}000 characters of each part and the storage audit covers 33 storage units.

\begin{table}[t]
\centering
\small
\setlength{\tabcolsep}{4pt}
\caption{Acquisition and parsing audit over the 100 submitted arXiv papers.
Values are asserted by the run scripts from saved artifacts, not hand-counted.
The content audit samples the first 1{,}000 characters of each part. Duplicates
are compared by fingerprint at a 0.90 threshold~\cite{manku-near-dup-2007}.}
\label{tab:acq-pipeline}
\begin{tabular}{@{}lr@{}}
\toprule
Check & Result \\
\midrule
Papers parsed and content-routed & 100/100 \\
Content-routed libraries & 5 \\
Stored parts after 30{,}000-character split & 322 \\
Per-paper ingest checks (nine checks) & 100/100 pass \\
True duplicates (35 near-duplicate pairs) & 0 \\
Descriptions upgraded & 165/165 \\
Content audit & 322/322 \\
Storage audit & 33/33 \\
Retrieval regression probes & 15/15 \\
Dense replica index (800-character chunks) & 4{,}528 chunks \\
\bottomrule
\end{tabular}
\end{table}

\textbf{Accounting gap.} This campaign does not support the claim that all 100 papers became usable, and we do not make it: 100 papers were submitted, the state file \texttt{exp10\allowbreak\_ingest\allowbreak\_state.json} records 91 names, and 79 were retrievable at evaluation time, so 21 papers were never retrievable. The reported failure modes are formula-dense scanned papers, which the parser handled poorly, and one crash of the tool layer during batch parsing that was restarted from a checkpoint. No artifact records a count of attempted against failed papers, and the state file has no per-paper parse status, so we cannot say whether the nine papers missing between submission and the recorded list failed or were never recorded. The absent counter is an observability gap: reconciliation is possible only against the artifacts that survive~\cite{missier-prov-2013}.

\textbf{Cross-lingual ingestion.} The benchmark ran twice over 400 pages per language. The ingest success rate was 1.0 in both rounds for all three languages. Membership accuracy, whether a stored page landed in the intended library, was 0.875 for English, 1.0 for Chinese and 1.0 for Japanese (400 pages per language). Storage completeness, computed on a sample of 40 pages per language with the shards of large documents summed, was 0.975, 0.9559 and 0.9402 respectively. Self-retrieval Hit@1, again at 40 pages per language, was 0.975, 1.0 and 1.0. Round~2 was bit-identical to round~1. The English figure of 0.875 is the weakest cross-lingual number we report, and it is an organisation defect rather than a parsing defect: the page was stored, but in the wrong library roughly one time in eight.

\textbf{Smaller campaigns.} A 50-paper, 26-domain campaign parsed and verified 50/50 papers, took 1{,}115.5\,s at 16.3--40.4\,s per paper, produced 202 tags, batch-indexed 50 documents and reported \texttt{all\_probes\_ok}; in its routing run 43 of the 50 documents carried \texttt{tag\_ok:false} while \texttt{verify\_ok} passed for all 50, so the tag-quality check failed often without blocking ingestion. Its 30{,}000-character tier turned 50 papers into 165 documents (nine whole documents and 156 parts), passed the storage check 50/50 and the description-and-tag check 165/165 with zero failures, stored 3{,}798{,}446 characters and took 2{,}672\,s to parse. Two caveats attach: one stage re-ran because of inbox ordering, and more than 800 extracted images were not uploaded. Descriptions and tags are checked against document content before storage~\cite{park-metadata-quality-2009}.

\subsection{E2: Organization, consistency, and management operations}
\label{sec:e2}

We define organisational correctness as three-source consistency: the files on disk, the tree index and the per-library metadata must agree about which documents exist, where they live and how much indexed content each has. Under that definition we found a real defect in a disposable test library. After a document was renamed, the vector index retained the chunks of the old name: the recorded chunk count was 18 where 9 were expected, a duplicate search matched the document against its own former path at 0.9738 similarity, and a vector probe returned every chunk twice with identical scores. Neither the automatic re-indexing that follows a rename nor a forced batch index removed the stale chunks; only an explicit force re-index repaired the document, after which the recorded count was 9. The failure is invisible from the metadata alone.

\textbf{Reorganisation campaign.} The full reorganisation workflow then ran over the 322-document arXiv corpus: 322/322 metadata checks and 165/165 description checks passed, the tag registry was de-duplicated, all 322 documents were batch-indexed, 15/15 retrieval regression probes passed with zero degradation, and 35 near-duplicate candidate groups (similarities 0.90--0.975) were adjudicated as zero true duplicates, every group being a set of adjacent parts of the same paper; no document was deleted (Table~\ref{tab:org-campaign}). Two honesty notes apply: the three-source spot check sampled 33 of the 322 documents, so it is a sample rather than a full pass, and sub-library splitting and hierarchy creation were analysed but deliberately not executed, pending user confirmation, so the campaign is not evidence for hierarchical restructuring.

\begin{table}[t]
\centering
\small
\setlength{\tabcolsep}{4pt}
\caption{Reorganisation campaign over the 322-document arXiv corpus.
The three-source spot check samples 33 documents; the description, tag and
duplicate checks cover the full corpus.}
\label{tab:org-campaign}
\begin{tabular}{@{}lrr@{}}
\toprule
Check & Scale & Result \\
\midrule
Metadata consistency (three sources) & 322 documents & 322/322 pass \\
Description check & 165 documents & 165/165 pass \\
Tag de-duplication & 4 groups & 4/4 pass \\
Batch index & 322 documents & 322 indexed \\
Duplicate adjudication & 35 groups & 0 true duplicates \\
Three-source spot check & 33 of 322 documents & 33/33 pass \\
Retrieval regression probes & 15 probes & 15/15, no degradation \\
\bottomrule
\end{tabular}
\end{table}

\textbf{Scale.} Growing the corpus from 50 to 100 papers (163 to 322 documents) did not cost retrieval quality: dense top-1 hits went from 48/50 to 49/50 and MRR from 0.98 to 1.00, and reorganisation itself produced a 2 percentage point improvement. The judgement field recorded with that comparison reads ``held (within $\pm$5\,pp)''; we report the field rather than restating the gain as a retrieval improvement.

\textbf{Lifecycle operations.} Table~\ref{tab:org-ops} lists the operations exercised. Create, move, delete and forced re-index behaved as specified: the move left the source library at zero documents and the target at one, and batch deletion removed both disposable test libraries. Two failure modes are worth reporting. The task-status tool reports a task as \texttt{running} after it has finished: the record carries a \texttt{finished\_at} timestamp while the status field still reads \texttt{running}, and the same symptom appeared in the reorganisation graph build, where node reconciliation replaced task polling. The update check timed out after 30\,s (\texttt{ragctl version -{}-json timed out after 30s}, success false) while every other probe in the batch passed; the project-status tool also silently ignores its version flag. Merge and split were never exercised under controlled conditions and we claim nothing about them. These 16 skill scenarios ran against two disposable test libraries in a standalone skill environment; the ingest tools were exercised in E1, not here.

\begin{table}[t]
\centering
\small
\setlength{\tabcolsep}{4pt}
\caption{Lifecycle operations exercised on disposable test libraries.
Merge and split were analysed but never executed.}
\label{tab:org-ops}
\begin{tabular}{@{}ll@{}}
\toprule
Operation & Observed outcome \\
\midrule
Create & Pass, document indexed and retrievable \\
Rename & Pass, stale chunks remained until forced re-index (18 recorded, 9 expected) \\
Move & Pass, source 1$\rightarrow$0$\rightarrow$1, target 0$\rightarrow$1$\rightarrow$0 \\
Delete & Pass, batch delete 1/1; both test libraries removed \\
Re-index & Automatic and batch paths did not purge; pass with \texttt{force=true} (18$\rightarrow$9) \\
Merge, split & Not evaluated \\
Task status & Fail, reports \texttt{running} after \texttt{finished\_at} is set \\
Update check & Fail, \texttt{ragctl version} timed out after 30\,s \\
\bottomrule
\end{tabular}
\end{table}

\textbf{Tags and graph.} The live tag registry held 316 tags, and the tag-quality check passed 165 documents with zero failures in 47.6\,s. Orphan tags do accumulate, among them test tags and section titles extracted from headings; the registry offers a dry-run detector and a cleanup operation, and in one case tags from a test library leaked into the production registry, the failure mode that vocabulary-based organisation is prone to~\cite{golder-tagging-2006}. The document graph is used for navigation and candidate expansion, never as evidence. Writing one document's four tag edges moved the edge count from 4{,}675 to 4{,}679, against 658 nodes, 333 documents and 316 tags; graph statistics for the 165-document corpus were 335 nodes and 1{,}671 edges.

%% n8b_eval.tex -- evaluation, middle part: E3 (operational knowledge and
%% automated distillation) and E4 (retrieval quality and the role of the
%% verdict). Assembled 2026-10-12. Labels owned here: sec:e3, sec:e4,
%% tab:retrieval, tab:e4-modes.
%% Numbers: restructure-plan-2026-10-12.md (\S C.1 items 8.4/8.5, \S D-3, \S H),
%% functional-test-evidence-2026-10-12.md, v2/s7_OLDresults.tex, and the
%% artifacts named inline (paper-kbs/experiments/results/, benchmark-suite/
%% results/, review/fulltest-20261012, review-2026-10-08.md).
\subsection{E3: Operational knowledge and automated distillation}
\label{sec:e3}

The experience store holds operational lessons as first-class retrievable
documents: an entry is addressed, indexed and searched through the same
interfaces as an ordinary document, and it carries a credibility tier that
decides whether it may be shown at all
(Section~\ref{sec:org-experience}, Table~\ref{tab:exp-tiers}). The design
claim is that a lesson about practice should be governed like any other
content object. The evidence collected for that claim is case evidence and we
report it as such.

Three measurements exist. (i)~One extraction campaign over a single small
knowledge base produced thirteen candidate drafts
(\texttt{benchmark-suite/results/experience\_smoke.json}:
\texttt{extract\_candidates}~13, \texttt{drafts\_created}~13). (ii)~One
end-to-end scenario was sedimented into one experience
(\texttt{exp-d6f4d47ae3d6}) and a later global experience search returned that
entry first, with a local similarity of 0.711
(\texttt{review/fulltest-20261012}, scenario A6). (iii)~The draft-approval
path was never exercised: the draft pool was empty when it was queried
(\texttt{experience\_drafts\_list} returned \texttt{count}~0 with no drafts),
so \texttt{experience\_draft\_read} and \texttt{experience\_draft\_approve}
could not run at all (\texttt{benchmark-suite/results/SKILL-E2E-REPORT.md},
lines~27 and~103). The heuristic extraction path was exercised only in
dry-run mode, where it returned a raw document dump that the quality gate
would reject (line~100).

One scenario and thirteen candidates cannot support an effectiveness claim,
and we do not make one. The larger part of the automated-distillation
pipeline is unverified: scheduled distillation, review and approval at scale,
application feedback, and credibility decay were not exercised end to end.
What the measurements show is narrower and sufficient for the design
argument: an operational lesson can travel the document path --- created,
indexed, returned first by search, with its tier attached --- and nothing more.

\subsection{E4: Retrieval quality and the role of the verdict}
\label{sec:e4}
\label{sec:eval-retrieval}

\begin{table}[!htbp]
\centering
\small
\setlength{\tabcolsep}{3pt}
\caption{Tool-layer retrieval on the 42-document deployed industrial corpus
(mean over 16 gold-annotated questions; top-5 document lists; latency per
call, warmed).}
\label{tab:retrieval}
\begin{tabular}{lccccc}
\toprule
Method & Hit@5 & Recall@5 & MRR@5 & Latency (mean) & Latency (median) \\
\midrule
BM25                  & 1.00 & 1.00 & 0.958 & 0.9\,ms    & 1\,ms \\
Dense (BGE-M3)        & 0.94 & 0.94 & 0.856 & 18\,ms     & 17\,ms \\
RRF (BM25$\|$Dense)   & 1.00 & 1.00 & 0.927 & 11\,ms     & 9\,ms \\
Graph (Neo4j)         & 0.56 & 0.53 & 0.562 & 68\,ms     & 68\,ms \\
Two-stage             & 1.00 & 0.97 & 0.927 & 77\,ms     & 40\,ms \\
Fusion-only           & 0.94 & 0.94 & 0.856 & 62\,ms     & 66\,ms \\
QDCVR-Hybrid          & 0.44 & 0.41 & 0.226 & 542\,ms    & 582\,ms \\
\midrule
\multicolumn{6}{l}{\emph{Derived (offline, same logged hybrid pools):}}\\
\quad Judge-as-filter$^{\dagger}$ & 0.94 & 0.94 & 0.856 & \multicolumn{2}{l}{$^{\dagger}$542\,ms / 582\,ms} \\
\quad RRF-fuse & 0.88 & 0.88 & 0.606 & \multicolumn{2}{l}{(same single call)} \\
\bottomrule
\end{tabular}
\par\smallskip
\raggedright\small
$^{\dagger}$\,Derived rows replay the QDCVR-Hybrid run's own logged
12-candidate pools and judge scores, so they reuse that call's latency; the
judge executes once either way.
\end{table}

Table~\ref{tab:retrieval} compares nine configurations on the sixteen
gold-annotated questions of the deployed industrial corpus (42 documents;
\texttt{paper-kbs/experiments/results/exp1\_retrieval.json} and
\texttt{exp5\_judge\_filter.json}). The tool-layer arms retrieve and never
answer: their latencies are retrieval latencies, and, because no generation
follows, they carry no token cost. With sixteen questions the Wilson 95\%
intervals on the headline proportions are wide (judged hybrid $[0.23, 0.67]$;
judge-as-filter $[0.72, 0.99]$; BM25 and RRF $[0.81, 1.00]$), so the
decomposition below, not the point estimates, is the claim. A separate
six-method retrieval-mode ablation exists in the released logs but covers two
questions; it is a smoke check at $n = 2$ and no claim rests on it.

\textbf{On this corpus, lexical retrieval is strong.} BM25 \cite{bm25}
reaches Hit@5 1.00 (MRR 0.958) at about a millisecond per call; reciprocal
rank fusion \cite{rrf-cormack} matches it in Hit@5 and follows in MRR (0.927);
the two-stage BM25+graph-then-dense pipeline reaches Hit@5 1.00 with
Recall@5 0.97, one gold of a two-gold question falling outside the top-5.
Dense retrieval with BGE-M3 \cite{bge-m3} misses exactly one question (Q16),
where four near-duplicate sibling reports from one incident family crowd the
top-5 --- a confuser-density failure, not a vocabulary one. Graph search over
the document graph \cite{graphrag}, limited to node-name and tag matching, is
the weakest standalone signal (Hit@5 0.56, MRR 0.562). Nothing here is a
retrieval advance: the corpus is curated, the questions are written in its
vocabulary, three configurations tie at Hit@5 1.00, and the one that uses only
words leads on MRR.

\textbf{The central result is the verdict's role, not its existence.} The same
judged scores enter three configurations: not used (fusion-only, Hit@5 0.94),
used to re-rank candidates (QDCVR-Hybrid, 0.44), and used as a filter that
keeps survivors in their fusion order (judge-as-filter, 0.94). On the eight
questions where re-ranking and filtering differ (Q2, Q4, Q5, Q6, Q8, Q9, Q11,
Q12), filtering wins all eight and loses none, with eight further ties on
Hit@5, and it restores MRR to 0.856 (exact binomial sign test, two-sided
$p = 0.008$); the only ranking movement is Q7, where the gold slips from third
to fifth inside the top-5. The verdicts were not wrong about which documents
are evidence --- they were wrong as an ordering. This matches the
judge-miscalibration pattern reported for model-generated preferences
\cite{zheng-llmjudge,panickssery-selfpref}: an individually plausible verdict
can be ordinally unreliable, and the fix is the placement of the verdict, not
its existence or its strength. Its price is measurable: judging
twelve candidates adds well under a second with the resident worker, and the
hybrid call's latency surcharge over fusion-only is the observed difference
between 542\,ms and 62\,ms per call. Two caveats travel with the test. It is
post-hoc and descriptive --- one of nine configurations evaluated on the same
sixteen questions --- and the gate is a single locally fine-tuned checkpoint
(Section~\ref{sec:gate}), so every verifier finding in this paper inherits
that choice.

\textbf{A second corpus reproduces the baselines, not the collapse.} Ten
questions over the 79 retrievable papers of the 100-paper arXiv ingestion
(8{,}031 chunks; \texttt{paper-kbs/experiments/results/exp10\_papers\_%
toollayer.json}) are answered at Hit@5 0.9 by every configuration except graph
search, which reaches 0.2. One question (BQ03) is missed by all eight
configurations, and its gold paper is one of the 21 papers that never became
retrievable in that campaign: the single shared miss is an ingestion failure,
not a retrieval failure. We therefore report both denominators --- 0.9 over
the ten questions as asked, 1.00 over the nine answerable ones --- and say
which is which. The judged collapse of the industrial corpus does not
reproduce on this corpus, where the same-family confuser cluster that drives
it is absent; the graph configuration's 0.2 tracks absent relation tagging,
since the bulk import bypassed the tagging path.

\textbf{Thresholds and margins.} All 191 judged candidates from the sixteen
hybrid runs passed the 0.5 acceptance threshold, the lowest logged score being
0.589, so at the threshold layer the gate filtered nothing on this corpus
(\texttt{exp1\_retrieval.json}). The margin sweep that follows replays these
logged candidate pairs rather than re-retrieving them: with the top-5 floor
unchanged, a margin of 0.15 keeps the gold in 16 of the 17 in-pool gold
instances and 0.20 keeps it in all 17
(\texttt{paper-kbs/notes/review-2026-10-08.md}). A leave-one-out procedure that
selects the margin per question from the other fifteen reaches a gold survival
of 0.9375, its single miss being Q11, which the shipped 0.15 margin also
loses, while a fixed 0.20 floor reaches 1.00
(\texttt{exp6\_loo\_margin.json}). The shipped floor is thus an engineering
constant, not a fitted parameter, and the calibration that would justify it
has not been run.

\textbf{Rephrasing shows where the fragility sits.} On a grid of sixteen
paraphrase questions (\texttt{exp6\_paraphrase.json}) the judged configuration
falls to Hit@5 0.5 (MRR 0.189), while judge-as-filter and fusion-only hold at
0.9375 (MRR 0.833) and dense retrieval stays at 0.9375; BM25 falls from 1.00
to 0.875 (MRR 0.958 to 0.755). Both the lexical baseline and the judged
re-ranking lose ground when the wording changes, and the configuration that
refuses to let the verdict order the list does not. The gate's weakness is in
the ordering it induces, not in the verdicts it returns.

\textbf{End to end, the modes differ in checkability.} Ten fixed questions ran
through three modes that share one harness, one model and one neutral prompt
template and differ only in retrieval affordance: the platform's full read
surface, a bare agent with file tools over the exported papers, and dense
retrieval over an 800-character replica index of the corpus (4{,}528 chunks).
Table~\ref{tab:e4-modes} reports the outcome. All three modes produce answers,
and the bare agent grounds as many of them as the platform
(10/10 vs.\ 10/10), but its answers name files where the platform names
checkable addresses, and dense retrieval cites an address once in ten
questions. The modes differ mainly in checkability rather than in whether an
answer was produced. Three caveats belong to this comparison: both dense
grounding misses were tool-availability aborts; grading is conservative, and
word-boundary matching zeroes the bare agent's plural ``cathodes'' on BQ08; and
two platform cells re-ran once after a harness timeout, correctness calls are
the authors' trace readings, and cache state was uncontrolled.

\begin{table}[!htbp]
\centering
\small
\setlength{\tabcolsep}{4pt}
\caption{Three-mode end-to-end comparison on ten fixed questions over the
100-paper arXiv corpus. \emph{Grounded}: strict arXiv-id citation of the
designated paper. \emph{Key-fact}: at least 60\% of the fixed gold key facts
matched at word boundaries. \emph{Cited}: the answer names a part, section or
table. Wall clock and cost are per question.}
\label{tab:e4-modes}
\begin{tabular}{lccccc}
\toprule
Mode & Grounded & Key-fact & Cited & Wall clock & Cost \\
\midrule
Platform (read surface)   & 10/10 & 10/10 & 9/10 & 116\,s & \$0.55 \\
Bare agent (file tools)   & 10/10 & 8/10  & 6/10 & 48\,s  & \$0.63 \\
Dense retrieval (replica) & 8/10  & 6/10  & 1/10 & 41\,s  & \$0.28 \\
\bottomrule
\end{tabular}
\end{table}

E4 establishes one thing. Content verification placed as a filter rather than
as a ranker is the difference between a citation and a claim: every returned
document is judged, and the ordering the reader sees is not the verdict's.
Verification before answering is not itself new \cite{selfrag,crag}; what is
measured here is where the verdict belongs in the pipeline and what it costs
there. E4 does not establish a retrieval advance --- BM25 reaches Hit@5 1.00
on this corpus --- so recall is not the contribution, and no margin, threshold
or paraphrase variant tested repairs the ordering role of the verdicts. That
is the question the failure analysis of the verifier takes up.

\begin{figure}[t]
\centering
\includegraphics[width=\textwidth]{fig3_retrieval}
\caption{Tool-layer retrieval over the 42-document production corpus
(sixteen gold questions, top-5 document lists). Lexical retrieval is
strong on this corpus: BM25 reaches Hit@5 1.00. The judged hybrid loses
more than half of its top-5 hits (0.44); placing the same verdicts in a
filter role recovers the fusion result (0.94) because the verdicts were
not wrong about which documents are evidence, only about their order.}
\label{fig:retrieval}
\end{figure}

\begin{figure}[t]
\centering
\includegraphics[width=\textwidth]{fig6_margin_variants}
\caption{Margin and paraphrase variants over the same 191 logged judged
candidate pairs. (a) The relative cut replays the logged pairs with the
top-5 floor unchanged: the shipped 0.15 margin loses one of seventeen
in-pool gold instances, a 0.20 floor loses none, and a leave-one-out
selection over the other fifteen questions lands at 0.9375. (b) Under
sixteen paraphrases the judged re-ranking collapses to Hit@5 0.50
(MRR 0.189) while the filter repair and dense retrieval hold at 0.9375
and BM25 falls to 0.875 (MRR 0.755).}
\label{fig:marginvariants}
\end{figure}

% n8c_eval.tex — evaluation 8.6-8.8: E5 agent behaviour/abstention/cost,
% E6 verifier failure, E7 non-functional behaviour.
% Every figure is attributed to its artifact; corpora are never mixed.

\subsection{E5: End-to-end agent behaviour, abstention, and cost}
\label{sec:e5}
\label{sec:agentic}\label{sec:abstention}\label{sec:composition}\label{sec:eval-agent}
%% The four labels above are compatibility aliases: n2_related.tex:110,122,139
%% and n6_retrieval.tex:38,53-54 were written against the pre-restructure label
%% set. They all resolve to this subsection (agent behaviour, abstention,
%% latency composition).

\textbf{The agentic grid.}
Four arms --- the Search, Librarian and Hybrid lanes plus a bare model with
no retrieval tools at all --- answered four questions (Q2, Q4, Q6, Q7)
twice each through the public chat API, with one engine (\texttt{omp} over
GLM-5.3-Flash) and a 20-turn budget. Citation outcomes agree on all 16 of
16 arm--question cells across the two rounds: every grounded cell named
the gold document (12/12); the bare arm --- the same model, told to answer
from its own knowledge --- cited a source in 0 of 4 cells in either round.
The de-contaminated fact-marker check (delete every answer line naming a
gold document from the answer's first 600 characters, then re-test the
remainder) is 12/12 in round 1 and 11/12 in round 2; the exception,
Hybrid-Q2, keeps its markers only inside lines that also cite the gold ---
an artifact of the deletion rule, and its citation hit stands. Six of the
twelve grounded cells answer from libraries of one or two documents, where
hitting the gold is a low bar; the substantive evidence is the six cells
over the 31-document process library. Lane operating points are in
Table~\ref{tab:lanes}. Where the bare arm is wrong, it is wrong in the way
that motivates retrieval: on the lamination question it gave a plausible
textbook diagnosis (interfacial adhesion failure, moisture ingress)
missing the corpus's recorded root cause, thermocouple drift leaving the
encapsulant under-cured.

\textbf{Abstention is a property of the reading-and-answering layer, not of
a threshold.}
Three out-of-corpus questions produced three explicit not-found reports
(47.9, 73.3 and 113.5\,s wall clock; turns were not recorded for this
probe family), each enumerating the scope it searched (five libraries, 42
documents), declining to answer from memory, and citing no document as
evidence. The retrieval tools themselves never abstain: six graded probes
at the tool layer --- three partially covered, three absent --- retained
10 to 12 of their 12 candidates at every tier, with top scores up to
0.983 on absent probes and 0.965 on partial ones; the acceptance
threshold of 0.5 never fired as a refusal. Honesty is therefore enforced
where the answer is written --- in the lane's read-then-answer contract
\cite{selfrag,feng-abstain,yin-knowwhat} --- not by a similarity cut-off.

\begin{figure}[t]
\centering
\includegraphics[width=0.62\textwidth]{fig7_ood}
\caption{Out-of-corpus probes at the tool layer: for every probe the
judged hybrid keeps 9--11 of its 12 candidates and dense retrieval passes
chunks above the 0.35 threshold---the threshold alone never abstains. At
the agent layer, all six probe runs returned explicit not-found reports
with zero fabricated citations.}
\label{fig:ood}
\end{figure}

\textbf{Where the time goes.}
In the twelve grounded cells, machine-side retrieval and verification
stayed sub-second (BM25 0.9\,ms mean, dense 18\,ms, two-stage 77\,ms,
judged hybrid 542\,ms mean) against an end-to-end wall clock of
36--91\,s. The gap is model turns: 11.7, 12.8 and 15.3\,s per turn for the
Search, Librarian and Hybrid arms --- roughly
12--15\,s per turn, the turn-dominated profile reported for agentic
retrieval \cite{singh-agentic-rag}, and the reason the platform moves
orchestration server-side and caps turns. Results repeat; wall clock does
not: Search-Q2 took 70\,s over 6 turns in round 1 and 165\,s over 11 turns
in round 2, Librarian-Q4 84\,s over 6 turns and 92\,s over 11.

\textbf{Cost.}
The same task set on three engines separates by one to two orders of
magnitude at the same quality tier (Table~\ref{tab:e5-cost}): all 18
scenarios pass with every answer gold-hit and no fabricated citation. The
cheapest engine pays a 50--90\,s harness cold start per call, so its wall
clock is about twice its engine time; the Claude harness's wall clock is
approximately its engine time.

\begin{table}[!htbp]
\centering
\small
\setlength{\tabcolsep}{4pt}
\caption{Three engines on one task set (18 runs; artifact:
\texttt{review/fulltest-20261012/SUMMARY.md}). Cost is per run in USD as
metered by the harness.}
\label{tab:e5-cost}
\begin{tabular}{p{4.4cm}rrrp{2.4cm}}
\toprule
Engine (surface) & Runs & Wall (s) & Turns & Cost per run \\
\midrule
Agent surface (\texttt{/api/kb/agent/chat}) & 7 & 50.5--175.9 & 2--9 & \$0.9--1.8 \\
Claude harness (\texttt{engine=claude}) & 6 & 23.2--52.4 & 5--8 & \$0.11--0.36 \\
OMP harness (\texttt{engine=omp}) & 5 & 33.5--171.4 & 2--9 & \$0.03--0.07 \\
\bottomrule
\end{tabular}
\end{table}

\subsection{E6: When the verifier fails}
\label{sec:e6}
\label{sec:inflation}\label{sec:eval-verify}   %% compatibility aliases (n2_related.tex:110)

\textbf{Same-domain score inflation.}
The 191 candidates judged across the sixteen hybrid tool runs --- twelve
per run except Q14, whose merge produced eleven --- carry no usable
separation between gold and non-gold. Every candidate scored at least
0.589, far above the 0.5 acceptance threshold, so on this corpus the
threshold rejects nothing. The 17 in-pool gold instances (Q16's gold never
entered its pool) have a median of 0.88 and the 174 non-gold candidates
0.86, a difference a two-sided Mann--Whitney test cannot distinguish from
chance ($p = 0.26$). Only five of the 17 golds rank in the top three of
the judged ordering and nine rank eighth or lower. Re-ranking with the
verdict therefore loses gold: the judged ordering reaches Hit@5 0.4375
(7/16; 0.44 in Table~\ref{tab:retrieval}), where the verdict used only as
a filter, or as one signal in rank fusion, holds 0.94
(Section~\ref{sec:e4}).

\textbf{The verdict is a model \cite{laya}, and a second model fails the
same way.}
We replayed the same 191 logged (query, candidate) pairs --- a replay, not
a re-retrieval --- through a generic natural-language-inference
checkpoint: 330M parameters, one question-as-context template, premises
truncated to roughly the first 500 characters after tokenization, with 59
of the 191 pairs hitting the cap. It scores below the verdict (Hit@5
0.3125, MRR 0.1375) and its separation is inverted: gold median 0.0113
against non-gold 0.05285 (Mann--Whitney $p = 0.120$). Truncation is not
the explanation, because six of the ten gold documents that fit the
context window whole still score below 0.01; nor is it an out-of-domain
effect, as unrelated topics still score corpus documents at up to 0.9973,
so an NLI threshold would abstain on nothing either. A span-level replay
inverts harder and significantly: over the 191 pairs and 3{,}489 span
inferences at roughly 320 characters per span, the max-pooled score is
0.4728 for gold against 0.91855 for non-gold ($p = 0.0074$), with Hit@5
0.1875. Rank fusion of the two signals repairs neither (Hit@5 0.3125, MRR
0.111458). Replacing the judge checkpoint is not the repair.

\textbf{Paraphrase.}
The verdict's agreement with gold collapses when the same questions are
rephrased; that grid is reported in Section~\ref{sec:e4} and not repeated
here.

\textbf{What the failure tracks.}
Neither a generic NLI checkpoint, nor a span-level variant of it, nor rank
fusion over both reaches the verdict's own quality; the second corpus
shows the same gate losing nothing there (Section~\ref{sec:e4}), so the
collapse tracks confuser density within one domain --- many
near-duplicate documents competing for the same question --- and not the
checkpoint, the domain, the language, or the corpus size
\cite{zheng-llmjudge,panickssery-selfpref,lightman-verify,yin-knowwhat}.
The verdict earns its place as a filter and loses it as a ranker.

\begin{figure}[t]
\centering
\includegraphics[width=\textwidth]{fig5_e2e}
\caption{Agentic runs. (a)~Wall time per run; printed numbers are agent
turns. (b)~Citation gold-hit rate vs.\ mean agent turns (horizontal offsets added for
readability): the hybrid lane reaches fully grounded answers with the
fewest turns, the librarian lane spends the most turns for exhaustive
recall, and the bare agent cannot ground citations at all.}
\label{fig:e2e}
\end{figure}

\begin{figure}[t]
\centering
\includegraphics[width=0.5\textwidth]{fig4_inflation}
\caption{Evidence-criterion judge scores from the sixteen QDCVR-Hybrid tool
runs: all 191 judged candidates. Gold documents (green, $n = 17$; Q16's
gold never entered its candidate pool) and same-domain non-gold survivors
(orange, $n = 174$) overlap almost completely: medians 0.876 versus 0.863,
Mann--Whitney $p = 0.26$. The 0.5 acceptance threshold separates both
groups from zero, not from each other.}
\label{fig:inflation}
\end{figure}

\begin{figure}[t]
\centering
\includegraphics[width=\textwidth]{fig8_anatomy}
\caption{Anatomy of one real query (Q2; single logged QDCVR-Hybrid call).
Each line is one of the 12 fused candidates; the green band is the top-5
window the answer agent reads in full. (a)~The shipped re-ranking: the
judge scores five extruder-PLC telemetry reports $0.887$--$0.926$ and the
gold root-cause report $0.846$, demoting it from fusion rank~1 to
rank~8 --- outside the read window. (b)~The same call with the verdicts
used only as a filter: all 12 pass, the fusion order stands, and the gold
is the first document read. Same pool, same scores, same verdicts; only
the role of the induced ordering differs.}
\label{fig:anatomy}
\end{figure}

\subsection{E7: Non-functional behaviour}
\label{sec:e7}
\label{sec:robustness}\label{sec:eval-nonfunctional}

\textbf{Stability.}
A 27-run suite --- three task types $\times$ three lanes $\times$ three
rounds on the live multi-base instance
(\texttt{review/stability-multiround-20260930}) --- is summarised in
Table~\ref{tab:e7-stability}. Session-trace replay confirms 27/27
lane-contract compliance, with per-cell median wall clock spanning
34--76\,s. Caveats: correctness calls are the authors' trace readings,
cache state was uncontrolled, two platform cells re-ran once after a
timeout, and the $\pm$35\,s variance sits in the answer-generation
window.

\begin{table}[!htbp]
\centering
\small
\setlength{\tabcolsep}{5pt}
\caption{The 27-run stability suite (three task types $\times$ three lanes
$\times$ three rounds).}
\label{tab:e7-stability}
\begin{tabular}{p{2.4cm}p{6.4cm}p{4.3cm}}
\toprule
Task type & Expected behaviour & Outcome \\
\midrule
Answerable & correct attention formula & 9/9 correct \\
Not covered & report that the library holds no such content & 9/9 correct not-found \\
Out of corpus & refusal (the librarian lane stops at the catalog layer) & 9/9 refusals \\
\bottomrule
\end{tabular}
\end{table}

\begin{figure}[!htbp]
\centering
\includegraphics[width=\textwidth]{fig9_robustness}
\caption{The robustness tier at a glance. (a)~Paraphrase grid: per-method
Hit@5 on the original questions (blue) and their paraphrases (orange);
BM25's ceiling is vocabulary-dependent while dense, fusion, and the
filter repair hold. (b)~NLI separation over the 191 logged candidates
(span scores from 3{,}489 span inferences): both granularities
\emph{anti-separate} --- gold medians sit below non-gold medians.
(c)~Domain, language, and scale shift to 79 ingested arXiv papers: the
judged collapse
disappears when no same-family confuser clusters exist.
(d)~Graded abstention probes: the threshold keeps 10--12 of 12
candidates in every tier.}
\label{fig:robustness}
\end{figure}

\begin{table}[!htbp]
\centering
\footnotesize
\setlength{\tabcolsep}{2.5pt}
\caption{Robustness checks: the question each probe was designed to answer
and the measured result. All numbers come from released logs.}
\label{tab:robustness}
\begin{tabular}{p{2.7cm}p{3.9cm}p{5.9cm}}
\toprule
Check & Question & Result \\
\midrule
Cut-margin LOO & Is the 0.20 margin an artifact of selecting on all 16
questions? & The smallest-margin rule chosen on 15 questions transfers to
15/16 holdouts (miss = Q11, which the shipped 0.15 also loses); the
deployed 0.20 retains all 16 holdout golds. The $+0.02$ pad is
engineering judgment, not an LOO value. \\
Paraphrase grid ($n{=}16$) & Does the tool layer depend on the questions'
vocabulary? & BM25 1.00$\to$0.88 (MRR 0.958$\to$0.755); dense 0.94;
judged 0.50; filter = fusion 0.94. Pooled over 32 variants: 15 wins,
0 losses, 17 ties ($p < 0.0001$). Failure and repair are wording-robust. \\
Second signal (NLI) & Is inflation specific to the Laya checkpoint? &
A generic Chinese NLI ranks \emph{worse} (Hit@5 0.31, MRR 0.138;
$\rho = -0.13$) and \emph{anti-separates} (gold median 0.011 \emph{below}
non-gold 0.053, $p = 0.12$); out-of-corpus topics still score corpus
documents up to 0.997. Naive span-max NLI is worse still (0.19;
anti-separation $p = 0.007$); the rank-fusion ensemble recovers nothing
(0.31). \\
Symbolic check (Q6) & Can rule checks cover what model judges miss? &
The watch value ``31.35'' appears in the gold (1/1) and in no confuser
(0/11): the numeric case the judge mishandled is solved by a trivial
substring rule. \\
Graded probes ($n{=}6$) & Does the threshold behave differently between
partial and absent coverage? & No: the threshold keeps 10--12 of 12
candidates (up to 0.983) in every tier. Agent layer: scoped disclosures
on partial probes (one reclassified as covered --- the corpus held a
verified rework case); 3/3 absent probes returned explicit not-found. \\
Second corpus (100 arXiv papers) & Does the judged collapse survive a
domain, language, and scale shift with no confuser clusters? & No
collapse at 79 indexed papers (8{,}031 chunks): judged = fusion =
filter = 0.90 (one shared miss --- that gold paper failed parsing).
The loss tracks confuser density, not domain, language, or size;
graph 0.20 without relation tagging (bulk import bypasses it). \\
\bottomrule
\end{tabular}
\end{table}

\textbf{Concurrency and integrity.}
A separate load campaign (\texttt{review/stress-20260923/STRESS-REPORT.md},
2026-09-23) wrote one library concurrently: 20 simultaneous document
creations finished in 0.8\,s, 20/20 HTTP 200, with no corruption of
\texttt{.tree-fs.json} or \texttt{.knowledge-base.yml} (list count 20);
20 simultaneous index builds finished in 4.1\,s (20/20, 620 chunks); 10
simultaneous experience creations finished 10/10. Lexical re-verification
found 3/3 probe documents and a vector top-3 check found one; the artifact
attributes that to homogeneous templates diluting the semantic signal, not
to data loss. This bounds concurrent writes only: no crash-mid-write or
partial-index adversarial test was run, so durability under interruption
remains untested.

\textbf{Integrity incidents.}
Two index-integrity incidents surfaced during evaluation setup --- one
document missing from its vector collection, one library whose collection
had been silently lost. Both were caught by the pre-measurement checks and
repaired through the platform's own path (\texttt{index-document},
\texttt{reindex}); the detection path, not the incident, is the point.

\textbf{Deployment reality.}
The live snapshot of 2026-10-12 shows a running system with 19 libraries,
including the 322-document arXiv corpus content-routed into five domain
libraries. Every live-system figure here is a timestamped snapshot of a
system in use, not a frozen benchmark; the five industrial libraries
behind the measurements above were then absent from the live catalog.

\section{Discussion and lessons learned}
\label{sec:discussion}

The evidence in Section~\ref{sec:eval} supports a distinction that the
retrieval literature tends to blur: finding a document and being allowed to
cite it are different operations, and the second one is a system property. The
lessons below are ordered from the retrieval layer outwards, because that is
the order in which we learned them.

\noindent\textbf{L1. Vectors are fast; content is accurate, and verification
has two costs.} The tool-layer comparison is the cleanest statement of the
retrieval subproblem: on a curated corpus a millisecond-level BM25 beats an
18\,ms dense retriever, and neither knows whether a document \emph{answers}
the question. The gate converts ranked candidates into accountable evidence
with per-citation provenance. Its costs are temporal (about 0.5\,s per dozen
candidates end to end; 1--2\,s per six-document batch for the resident worker)
and, on this corpus, accuracy-related: 0.94\,$\to$\,0.44 Hit@5 under
same-domain inflation when the verdicts are used as an ordering. Auditability
is what the gate buys today; ranking quality is what calibration must earn
back.

\noindent\textbf{L2. Latency engineering is turn engineering.} With retrieval
sub-second, the only lever that matters for interactive latency is the number
of agent turns. Tool whitelists that remove detour tools, one-call server-side
pipelines, and small judge batches that avoid transport timeouts together took
the hybrid lane from about 29 turns in the earliest serial-skill configuration
to 3--5 turns in the present measurements. The same accounting explains the
engine comparison: the third engine's wall clock exceeded its engine-side time
by 50--90\,s per call, which is harness startup, not retrieval.

\noindent\textbf{L3. The verifier is a model too, so calibrate it.} The
inflation analysis is the paper's cautionary result: a verifier inherits
domain-internal blind spots, and a threshold calibrated on cross-domain
negatives (0.5) does not separate same-domain near-misses. Fail-closed
semantics and forced reading contain the damage; they do not replace
calibration. The second-signal probes sharpen the point: a generic entailment
model fails too, but differently, anti-separating golds while still rejecting
no out-of-corpus topic, and a symbolic rule decides the numeric case that the
shipped judge mishandles. Verifier choice changes which errors you get, not
whether you get any.

\noindent\textbf{L4. Honesty needs an enforcement layer, not a prompt.} Lane
separation works because the server physically withholds tools: a
librarian-lane agent cannot call the vector tool even if instructed to. The
out-of-corpus probes show the same asymmetry for abstention. A threshold
cannot refuse; it accepted every out-of-corpus candidate it was shown. The
reading-and-answer layer refused all six probe prompts with disclosed searches
and zero fabricated citations. A bare agent without the platform refused
explicitly in two of four runs, and in the others produced confident,
unsourced, partly wrong answers whose wrongness no audit trail exposes.
Unassisted honesty is a coin flip; enforced honesty is inspectable.

\noindent\textbf{L5. Production incidents taught honest observability.} Three
early incidents shaped the implementation: a corrupted index segment that made
one library silently un-retrievable across lanes, an out-of-memory indexing
run that returned success with an empty payload, and judge failures that had
to reject candidates rather than pass them. The common thread is that every
layer reports what it actually verified and every consumer treats the
unverifiable as failed. Table~\ref{tab:failures} consolidates the incident log
into the failure-mode catalogue we wish we had on day one, extended with the
three defects the organization and functional campaigns surfaced.

\begin{table}[!htbp]
\centering
\footnotesize
\setlength{\tabcolsep}{4pt}
\caption{Failure-mode catalogue from production and campaign use: how each
failure presents, the detection signal, the fix, and the residual risk an
adopter still carries.}
\label{tab:failures}
\begin{tabular}{p{2.9cm}p{3.1cm}p{3.0cm}p{3.1cm}}
\toprule
Failure & Presents as & Detection signal & Fix / residual risk \\
\midrule
Corrupted index segment (silent loss) & ``Success, zero results'' for one
library across lanes & Treat zero-result success as suspect; cross-lane
census & Re-index the collection; residual: any new silent-loss mode needs
the same tripwire \\
Out-of-memory during indexing & Success envelopes with empty vector payloads
& Judge success by chunk counts, not envelopes & Memory hygiene between
embedding and judging workloads; residual: other silent truncations \\
Driver contention (Windows) & Embedding latency 68\,s $\to$ 27\,min &
Latency budget on the embedding path & Isolate embedding and judging
workloads; on Linux, watch VRAM contention instead \\
Judge worker cold start & First judge batch stalls 10--12\,s & Worker
residency heartbeat & Keep the resident worker warm; residual: first-call
timeouts still fail closed \\
Un-whitelisted harness tools & 14--170\,s detours per run & Turn count and
tool-call audit log & Server-side per-lane whitelists; residual: new harnesses
need new whitelists \\
Judge or engine errors & Wrong or missing verdicts & Fail-closed rejects at
the gate & Errors reject, never pass; residual: availability, not
correctness \\
Rename without re-index & Chunk count 18 where 9 exist; the document matches
its own former path & Three-source audit; near-duplicate report &
Force re-index; residual: the write path still cannot re-index implicitly \\
Graph node not created & A library reports two of three documents &
Per-library document count versus graph overview & Rebuild the graph for that
document; residual: construction is not yet guaranteed by ingestion \\
Task record outlives its task & A finished build still reports ``running'' &
Compare task status against the artifact it was to produce & Query the
artifact, not the task; residual: status is not completion evidence \\
\bottomrule
\end{tabular}
\end{table}

\noindent\textbf{L6. The acquisition gate sets the ceiling for everything
above it.} In the 100-document campaign, 21 documents failed to parse, mostly
dense scanned material with formulas and tables. The system's value here is
not that the rate is low, but that the failure is a reported outcome: a
parse failure that is stored as an empty or garbled document would enter the
catalog, receive a description, be indexed, and be retrieved and cited by the
same machinery that cites good documents, with no signal anywhere that
anything was lost. The same argument applies to a missing index or a stale
description. This is why the acquisition layer owns a post-condition rather
than a warning, and why the knowledge layer is a component of the system
rather than a precondition to retrieval.

\noindent\textbf{L7. Organization is a retrieval budget.} Content routing
decides what the candidate generator can see. A catalog that is curated and
consistent lets an agent choose a library and a document set from metadata
alone, which is the mechanism behind the catalog-first lane and behind
cross-library discovery in a deployment with nineteen libraries. An
ungoverned catalog produces the opposite: candidate generation degenerates
into a global vector search whose recall is limited by the cap, and no
downstream verification can recover a document that was never proposed. The
organization layer is therefore not administrative overhead around retrieval;
it is where retrieval's recall is decided.

The lessons share one implication for how systems of this kind should be
evaluated. A claim about enforcement is a claim about a control, so it needs a
matched ablation: the same pipeline with the mechanism disabled, the same
questions, the same corpus. Where we have that control (retrieval with and
without verification at the tool layer) the claim is quantified; where we did
not have it at first (an end-to-end agent arm with the gate disabled) we fixed
the protocol before running it rather than after
(Section~\ref{sec:eval-agent}). The verifier's negative result is the reason
we insist on the distinction: without the matched ablation, a filter and a
re-ranker are indistinguishable from the outside, and only one of them is
defensible.

\section{Threats to validity}
\label{sec:threats}

\noindent\textbf{Scale.} The retrieval corpus at the centre of the evaluation
is 42 documents across five libraries; the acquisition and organization
campaigns scale that to 322 documents, and the second corpus adds 79. Absolute
Hit@5 values (BM25 at 1.00 on the industrial corpus) will not transfer to
collections orders of magnitude larger, where the hybrid lane's candidate cap
and the librarian's shelf pruning become load-bearing rather than optional.
The benchmark suite shipped with the platform is the intended scaling path,
and a hundreds-of-documents run with an author-independent question set
remains the full answer.

\noindent\textbf{Question set.} Sixteen tool-layer and four agent-layer
questions on the industrial corpus, ten on the second corpus and forty cloze
questions per language in the acquisition campaign were authored by the
system's developers. Gold documents and fact markers are released for
independent scrutiny, but the samples are small and favour the corpora's
vocabulary. The paraphrase grid quantifies part of that exposure, since
lexical ceilings drop under rewording while the verifier findings hold. One
circularity is named rather than silently enjoyed: Q16's question embeds its
gold document's filename token, which guarantees lexical success on that
question; the fix is scheduled, and because Q16 is the one question on which
the fusion stage loses a gold, the correction can only tighten the reported
decomposition. The second corpus carries a related exposure that we disclose
instead of exploiting: the acquired arXiv documents are named
\texttt{\{domain\}\_\_\{identifier\}\_\_\{slug\}}, so content-routing accuracy
against the topical prefix is not an independent measurement of routing
quality, and we report routing coverage rather than routing accuracy
(Section~\ref{sec:e2}).

\noindent\textbf{Experimental controls.} The tool layer isolates the gate's
contribution cleanly: its nine configurations include retrieval without
verification and the same retrieval with the gate in either role, so the
ranking-layer claim rests on a matched control. The agent layer originally did
not carry the corresponding control, because its four arms are three grounded
lanes (retrieval \emph{and} verification) plus a bare agent (neither), which
cannot separate the contribution of retrieval from that of verification. The
missing arm, a grounded lane with the gate disabled and the fusion order
returned, was added at the end of this revision with its protocol fixed before
the run (Section~\ref{sec:eval-agent}), together with a second ablation that
varies the search partition rather than the gate
(Section~\ref{sec:e2}). Both inherit the same questions, libraries,
engine and turn budget; neither was tuned after its results were seen.

\noindent\textbf{Evidence provenance is uneven.} The evidence groups differ in
how they were produced, and the paper reports the difference rather than
smoothing it. The acquisition, organization and functional campaigns are
audited runs whose per-check outcomes are released. The experience subsystem,
by contrast, is supported by case evidence: one distillation campaign with
thirteen candidate drafts, a single sedimented experience, and an approval
path that was never exercised. The full 100-document ingestion exists as a
derived state file rather than as a timed, per-document log, and its
accounting is not fully closed: 100 documents were submitted, 91 names appear
in the derived state, and 79 documents are present in the library. Part of the
gap is the 21 reported parse failures, but the remainder is unattributed, so
the parse-failure rate should be read as a bounded estimate rather than as an
exact measurement. We state this because a system paper's most useful
contribution may be the honesty of its ledger.

\noindent\textbf{Judge and engine dependence.} Scores and turn counts come
from one judging checkpoint and one agent engine per campaign; lanes, not
models, are the object of comparison, but absolute numbers will move with
engine versions. The deployed gate remains a single checkpoint, although the
two strongest verifier findings are no longer single-signal: a generic
entailment model fails too, but differently (it anti-separates golds and still
rejects no out-of-corpus topic), and a symbolic rule decides the numeric case
the shipped judge mishandles. Trained evidence-verification models and learned
signal ensembles remain open; the naive span-max probe and the rank-fusion
ensemble answer the two easiest upgrades negatively.

\noindent\textbf{Design surface.} Three lanes triple the policy surface an
adopter must test, and the reported evidence supports two distinct operating
points, with the catalog-first lane as the exhaustive option and the
orchestrated lane as the latency-tuned interactive default. A two-lane
deployment is defensible today, and query-shape auto-selection is future work.

\noindent\textbf{Consumer-side evidence.} The honesty contract is evaluated
from the system's side, that is, that not-found reports are well formed and
non-fabricated, not from the reader's side. Whether an operator acts on a
not-found report, escalates to a human, or ingests the recommended source is
not measured, and this is the paper's main missing external-validity leg. The
natural study is a small operator trial over a fixed probe set with two
conditions, platform reports versus the same model's unaided answer, scored on
the downstream action taken; the probe set and transcripts are released so it
can be run independently.

\noindent\textbf{Nondeterminism.} Agent runs are stochastic. We report two
full repetitions of the sixteen-cell grid (citation agreement 16/16, with wall
time and turn counts fluctuating substantially), retry only zero-turn
transport flakes, and release all per-run logs including timeout re-runs.
Fact-marker checks are automated substring matches; we additionally report
decontaminated checks computed after removing answer lines that name a gold
document, and the out-of-corpus answers were audited manually. Deterministic
layers were verified to be bit-identical across two rounds except for latency
fields.

\section{Conclusion}
\label{sec:conclusion}

QDCVR turns ``the retriever found it'' into ``the system verified it''. Its
design position is that the governed object of a knowledge-base management
system is content: content decides where a document belongs, content decides
whether a candidate may be cited, and the absence of content that answers the
question is itself an outcome the system must report. That position is
implemented as contracts rather than conventions (RQ1): parse gates that
reject visibly, placement rules that read the document, a three-source
consistency audit over the disk, the tree registry and the library metadata,
lane policies enforced by a server-side tool whitelist, and verdicts that fail
closed.

The five research questions are answered by measurement rather than
assertion. Acquisition and content routing hold at the scale we tested, while
21\% of submitted papers failed to parse and were rejected rather than stored
(RQ2). Staged, multi-channel retrieval with content adjudication reaches
documents that one-shot vector top-$k$ misses, and the benefit comes from
staged recall plus content verification, not from the fusion of the two
retrieval stages (RQ3). Verification costs well under a second and buys
per-citation provenance, but as an ordering it inherits a same-domain blind
spot that we quantify and partially repair; the failure tracks confuser
density rather than the checkpoint, domain, language or scale (RQ4). Honest
abstention lives in the reading and answering layer, never in the threshold,
end-to-end latency belongs to model turns, and engine choice moves cost by two
orders of magnitude at the same quality tier (RQ5).

The roadmap follows from the same ledger. Verifier calibration against
per-library null distributions, trained evidence-verification models, and the
verdict-as-filter role that the tool layer already supports are the immediate
next steps; larger-corpus evaluation with an author-independent question set,
consumer-side evidence that not-found reports serve their readers, and lane
auto-selection from query shape are the open empirical ones. The direction is
towards retrieval systems whose trustworthiness is an inspectable property of
the pipeline rather than a hope about the model.

\section*{Data availability}
The platform source code, the acquisition and organization campaign logs, the
tool-layer and agentic experiment scripts, the gold question set with fact
markers, the paraphrase grid, the leave-one-out margin replay, the
second-signal and span-level NLI replays, the graded abstention probes, the
cross-lingual acquisition results, the gate and partition ablation runs, and
all raw per-run JSON logs (tool-layer retrieval rows and agentic transcripts,
including the timeout re-runs) are available at
\url{https://github.com/kingdol666/rag-knowledge} (benchmark suite:
\texttt{benchmark-suite/}; paper artifacts and evidence inventories:
\texttt{paper-kbs/}).

\section*{Ethics statement}
This article does not report studies involving human participants, human
data, or animals.

\section*{Declaration of generative AI in the writing process}
During the preparation of this work the authors used LLM-based coding
agents to assist with experiment execution, data tabulation, figure
preparation, and manuscript drafting. All experimental numbers originate
from logged runs of the deployed system on the authors' machine; the
authors reviewed and take full responsibility for the content of the
published article.

\section*{Declaration of competing interest}
The authors declare that they have no known competing financial interests
or personal relationships that could have appeared to influence the work
reported in this paper.

\section*{Author contributions (CRediT)}
\textbf{Author 1:} Conceptualization, Methodology, Software, Investigation,
Data curation, Writing -- original draft. \textbf{Author 2:} Validation,
Writing -- review \& editing, Supervision. %% AUTHOR ACTION: adjust.

\section*{Acknowledgments}
This research did not receive any specific grant from funding agencies in
the public, commercial, or not-for-profit sectors. %% AUTHOR ACTION

\appendix
\setcounter{table}{0}
\section{Per-question tool-layer results}
\label{app:perq}
Table~\ref{tab:appendixranks} breaks the tool-layer comparison down to
individual questions: the best gold rank each configuration reached in
its top-5 window, for every gold-annotated question. Reading the columns
side by side localizes every aggregate number in the paper: BM25 and RRF
miss at most one question (Q16's confuser family and Q7's telemetry
cluster account for the remaining sub-1.0 MRRs); the judged hybrid's nine
misses are exactly the questions where same-domain confusers outranked
the gold (Q2, Q4, Q5, Q6, Q8, Q9, Q11, Q12) or crowded it out upstream
(Q16); and judge-as-filter restores the gold on every one of the eight
judging losses while inheriting Q16's upstream failure --- the
decomposition that Fig.~\ref{fig:marginvariants}(b) aggregates and
Fig.~\ref{fig:anatomy} dissects one query at a time.

\input{tab_appendix_ranks}


\bibliographystyle{unsrtnat}
\bibliography{refs,refs_new_verified}

\end{document}
